% Adaptación de FORMATOAPE.tex para APE3; original preservado.

\documentclass[11pt,a4paper]{article}
\usepackage[utf8]{inputenc}
\usepackage[T1]{fontenc}
\usepackage[spanish]{babel}
\usepackage{geometry}
\usepackage{graphicx}
\usepackage{array}
\usepackage{tabularx}
\usepackage{makecell}
\usepackage{hyperref}
\usepackage{xcolor}
\usepackage{fancyhdr}
\usepackage{titlesec}
\usepackage{helvet}
\usepackage{ragged2e}
\usepackage{enumitem}
\usepackage{setspace}
\usepackage{amsmath}
\usepackage{booktabs}
\usepackage{float}
\usepackage{listings}
\usepackage{hyperref}
\usepackage{xurl}
\usepackage{microtype}
\usepackage{float}
\emergencystretch=2em
\renewcommand{\familydefault}{\sfdefault}
\geometry{a4paper,left=3cm,right=3cm,top=2.35cm,bottom=2.35cm,headheight=1.9cm,headsep=0.45cm}
\setlength{\parindent}{0pt}
\setlength{\parskip}{0.3em}
\setlength{\tabcolsep}{5pt}
\renewcommand{\arraystretch}{1.13}
\pagestyle{fancy}
\fancyhf{}
\fancyhead[L]{\IfFileExists{reporte_tecnico_overleaf_logo.png}{\includegraphics[height=1.35cm]{reporte_tecnico_overleaf_logo.png}}{\small Universidad Nacional de Loja}}
\fancyhead[R]{\small FEIRNNR -- Carrera de Computación}
\fancyfoot[C]{\thepage}
\renewcommand{\headrulewidth}{0.5pt}
\renewcommand{\footrulewidth}{0pt}
\titleformat{\section}{\bfseries\fontsize{15}{17}\selectfont}{\thesection.}{0.8em}{}
\titlespacing*{\section}{0pt}{1.0ex}{0.3ex}
\newcolumntype{L}[1]{>{\RaggedRight\arraybackslash}p{#1}}
\newcolumntype{Y}{>{\RaggedRight\arraybackslash}X}
\lstset{basicstyle=\ttfamily\footnotesize,breaklines=true,frame=single,showstringspaces=false,columns=fullflexible}
\hypersetup{hidelinks}

% Soporte mínimo del contenido académico extenso.
\usepackage{longtable}
\usepackage{calc}
\usepackage{fvextra}
\providecommand{\tightlist}{\setlength{\itemsep}{0pt}\setlength{\parskip}{0pt}}

\begin{document}
\thispagestyle{fancy}
\begin{center}
{\bfseries\fontsize{21}{24}\selectfont Reporte Técnico de Actividades\\ Práctico-Experimentales Nro. 003\par}
\vspace{0.25cm}
{\large Arquitectura, OpenAPI y endpoints base.}
\end{center}

\section{Datos de Identificación del Estudiante y la Práctica}
\begin{tabularx}{\textwidth}{|L{0.43\textwidth}|Y|}
\hline
\textbf{Nombre del estudiante(s)} & [PENDIENTE: integrantes reales del equipo] \\ \hline
\textbf{Asignatura} & Desarrollo Basado en Plataformas \\ \hline
\textbf{Ciclo} & 5 A \\ \hline
\textbf{Unidad} & 1 \\ \hline
\textbf{Resultado de aprendizaje} & Discute cómo los estándares Web influyen en el desarrollo de software, bajo los principios de solidaridad, transparencia, responsabilidad y honestidad. \\ \hline
\textbf{Práctica Nro.} & 003 \\ \hline
\textbf{Título de la Práctica} & Arquitectura, OpenAPI y endpoints base. \\ \hline
\textbf{Nombre del Docente} & Edison Leonardo Coronel Romero \\ \hline
\textbf{Fecha} & Semana 3; [PENDIENTE: fecha real de la sesión] \\ \hline
\textbf{Horario} & Por definir según horario académico \\ \hline
\textbf{Lugar} & Aula / laboratorio asignado \\ \hline
\textbf{Tiempo planificado en el Sílabo} & 3 horas \\ \hline
\end{tabularx}
\textbf{Proyecto:} SIGEAC / SEC-UNL.\par
\textbf{Estado:} artefactos y pruebas parciales sobre Spring Boot; DRF, sobre universal, PostgreSQL y cierre docente pendientes.\par

\textbf{Revisión posterior:} consultar \path|REVISION_CUMPLIMIENTO.md|
para el dictamen por paso y los pendientes de la guía. Las verificaciones
documentales repetidas no son una nueva ejecución del backend.\par

\section{Objetivo(s) de la
Práctica}\label{informe-objetivos-de-la-pruxe1ctica}

\begin{itemize}
\tightlist
\item
  Diseñar y documentar la arquitectura inicial del backend a partir de
  los requisitos y procesos priorizados.
\item
  Representar la solución mediante C4 (C1, C2 y C3) y UML/PlantUML para
  los elementos de código relevantes.
\item
  Refinar el contrato de API y documentarlo mediante OpenAPI / Swagger.
\item
  Implementar endpoints base coherentes con los requisitos, criterios de
  aceptación y semántica HTTP.
\item
  Mantener trazabilidad entre especificación, arquitectura, contrato e
  implementación.
\end{itemize}

Los endpoints demostrados ya existían: se contrastaron y ejecutaron, sin
atribuirles implementación nueva en APE3. El paso de implementar con
Django REST Framework y el contrato universal de respuestas están
pendientes de decisión e implementación, o de excepción docente
explícita.

\section{Materiales, Reactivos, Equipos y
Herramientas}\label{informe-materiales-reactivos-equipos-y-herramientas}

\subsection{Insumos reales}\label{informe-insumos-reales}

Se revisaron
\path|documentos-academicos/APE2-semana-2/specification.md|, el dossier
y matriz de AA2, y los procesos, decisiones y contrato de ACD3. El
\href{../ACD3-semana-3/REGISTRO_OBSERVACIONES.md}{registro de ACD3}
contiene hallazgos de preparación, no retroalimentación de una sesión
efectivamente registrada. No se inventaron acuerdos docentes para
resolverlos.

Se inspeccionaron los controladores/servicios de catálogo, procesos,
identidad y modelos; la constitución; ADR existentes; configuración y
pruebas. \url{FUENTES.md} enumera las rutas y referencias primarias.

\textbf{Herramienta solicitada }: Git/GitHub y GitFlow

\textbf{Estado observado / uso en esta preparación }: Repositorio Git,
rama main; no se creó rama, commit, PR ni publicación. Se propone flujo
en NOTA\_INDIVIDUAL.md.

\medskip\hrule\medskip

\textbf{Herramienta solicitada }: GitHub Spec Kit

\textbf{Estado observado / uso en esta preparación }: Estructura
\path|.specify/|, constitución y plantillas existentes; no se ejecutó
una cadena de comandos Spec Kit ni se simula su resultado.

\medskip\hrule\medskip

\textbf{Herramienta solicitada }: Python/Django/DRF

\textbf{Estado observado / uso en esta preparación }: Python 3.14
disponible; Django/DRF no encontrados en el intérprete inspeccionado.
Python se usó como cliente HTTP, no como backend.

\medskip\hrule\medskip

\textbf{Herramienta solicitada }: Backend real

\textbf{Estado observado / uso en esta preparación }: Java 25 instalado;
pom configura Java 21 y Spring Boot 4.1.1; backend arrancado desde
clases actuales.

\medskip\hrule\medskip

\textbf{Herramienta solicitada }: PostgreSQL

\textbf{Estado observado / uso en esta preparación }: Persistencia
principal del producto; no se ejecutó una aceptación PostgreSQL en esta
actividad.

\medskip\hrule\medskip

\textbf{Herramienta solicitada }: SQLite

\textbf{Estado observado / uso en esta preparación }: Perfil auxiliar
existente \path|sqlite-test|, base aislada en backend/target para esta
verificación. No sustituye PostgreSQL.

\medskip\hrule\medskip

\textbf{Herramienta solicitada }: Structurizr DSL

\textbf{Estado observado / uso en esta preparación }: workspace.dsl
elaborado con vistas C1/C2/C3; parser/render oficial no disponible en
PATH y no ejecutado.

\medskip\hrule\medskip

\textbf{Herramienta solicitada }: PlantUML

\textbf{Estado observado / uso en esta preparación }: Fuentes de clases
y secuencia elaboradas; renderizador no disponible y no ejecutado.

\medskip\hrule\medskip

\textbf{Herramienta solicitada }: OpenAPI/Swagger

\textbf{Estado observado / uso en esta preparación }: Dos contratos
académicos, Redocly instalado y validación de muestras con Ajv;
disponibilidad HTTP de Swagger comprobada.

\medskip\hrule\medskip

\textbf{Herramienta solicitada }: Cliente HTTP

\textbf{Estado observado / uso en esta preparación }: Python urllib con
cookie jar y peticiones reales; Postman CLI disponible, pero no
utilizado en esta campaña.

\medskip\hrule\medskip

\textbf{Herramienta solicitada }: Quarto

\textbf{Estado observado / uso en esta preparación }: Fuente técnica QMD
y salida HTML del paquete; ver estado de render en evidencias.

\medskip\hrule\medskip

\textbf{Herramienta solicitada }: Editor / LaTeX

\textbf{Estado observado / uso en esta preparación }: Archivos
versionables editados; FORMATOAPE.tex adaptado. Sin motor TeX local, sin
PDF compilado.

\medskip\hrule\medskip

\section{Procedimiento / Metodología
Ejecutada}\label{informe-procedimiento--metodologuxeda-ejecutada}

\subsection{Inicio: alcance, requisitos y decisiones
previas}\label{informe-inicio-alcance-requisitos-y-decisiones-previas}

Se conserva el recorte propuesto en Semana 2: acceso, catálogo académico
y preparación/consulta de procesos. F-ACD3-01 mantiene carrera y su
historial; F-ACD3-02 prepara un proceso con carrera válida y versión
publicada. Se excluyen autorización institucional, activación,
evidencias, revisión, evaluación, informes y archivo, aunque existan en
el producto completo.

Actores: usuario habilitado A-AA-01; administrador funcional A-AA-02;
gestor autorizado A-AA-03; lector autorizado A-AA-04. Las capacidades
del código no acreditan por sí solas competencias institucionales. Las
necesidades N-AA-01--05 y criterios completos siguen en el dossier AA2.

\subsection{Paso 1: consolidar el análisis
textual}\label{informe-paso-1-consolidar-el-anuxe1lisis-textual}

\textbf{F-ACD3-01, mantener catálogo.} Participantes: administrador,
frontend, API, aplicación, repositorio y auditoría. Inicio: solicitar
alta o corrección de carrera. Actividades: capturar datos; validar DTO;
comprobar identidad/capacidad; consultar unidad activa; verificar código
único; comprobar revisión en cambios; registrar carrera/revisión;
conservar historial; registrar auditoría; devolver resultado.
Decisiones: alta o cambio, entrada válida, unidad activa, código
disponible, revisión vigente. Datos: código, nombre, unidad, vigencia,
estado, ID, actor, fecha, revisión y correlación. Excepciones:
400/401/403/404/409/412/422/428 según causa. Resultado: recurso
conservado o rechazo controlado; desactivar conserva historia y no
borra.

\textbf{F-ACD3-02, preparar proceso.} Participantes: gestor, cliente,
aplicación de procesos, catálogo de carreras/modelos, persistencia y
auditoría. Inicio: solicitar preparación para una carrera. Actividades:
seleccionar carrera visible; seleccionar versión exacta; capturar código
y período; comprobar permiso/ámbito; verificar actividad y vigencia;
resolver versión publicada; crear PREPARATION; guardar proceso y
auditoría; devolver referencia; consultar proceso. Decisiones: permiso,
carrera utilizable, período contenido, publicación y visibilidad. Datos:
careerId, code, name, periodStart, periodEnd y modelVersionId; servidor
asigna ID, revisión 0, actor y fecha. Resultado: preparación, sin
activación ni autorización automática.

La verificación usó una carrera nueva para catálogo y la carrera piloto
sintética ya habilitada para la preparación, pues \path|CAREER_ADMIN| no
concede \path|PROCESS_MANAGE| automáticamente sobre una carrera recién
creada. No se cambió la política de autorización para facilitar la
demostración. El análisis extendido permanece en
\href{../ACD3-semana-3/ACD3_CLINICA_ARQUITECTURA_API.md}{ACD3}, sección
2.

\subsection{Paso 2: arquitectura C4 y
UML}\label{informe-paso-2-arquitectura-c4-y-uml}

\begin{itemize}
\tightlist
\item
  C1: administrador, gestor y lector se relacionan con SIGEAC; no se
  dibujan APIs institucionales no confirmadas.
\item
  C2: frontend Angular, backend monolítico modular Spring Boot y
  PostgreSQL. SQLite pertenece exclusivamente al entorno auxiliar
  probado. El almacén de evidencias del producto completo queda fuera
  del recorte.
\item
  C3: identidad/acceso, API de catálogo/procesos, aplicación de
  catálogo, aplicación de procesos, dominio, puerto de modelos,
  persistencia y auditoría. Las llamadas entre módulos son internas al
  backend, no HTTP entre microservicios.
\end{itemize}

Artefacto: \href{arquitectura/workspace.dsl}{workspace.dsl}. UML:
\href{arquitectura/catalogo-clases.puml}{clases de catálogo} y
\href{arquitectura/alta-carrera-secuencia.puml}{secuencia de alta},
basados en nombres reales. Los diagramas no se presentan como
renderizados o validados por sus motores oficiales. La
\href{../ACD3-semana-3/ARQUITECTURA_C4.md}{vista C4 previa} conserva una
visualización Mermaid en Markdown.

\subsection{Paso 3: decisiones
arquitectónicas}\label{informe-paso-3-decisiones-arquitectuxf3nicas}

Se conservan ADR-0001 (monolito modular), ADR-0004 (sesiones/ámbito) y
ADR-0018 (compatibilidad). Se elaboran propuestas académicas
\href{adr/APE3-001-alineacion-tecnologica.md}{APE3-001}, tecnología
exigida frente a existente, y
\href{adr/APE3-002-contrato-universal.md}{APE3-002}, contrato común y
migración. Ninguna se marca aceptada ni modifica los ADR del producto.

La pregunta sobre Spring/DRF se planteó al usuario durante esta
preparación. Sin elección expresa se conserva el producto y se prepara
la entrega sobre su implementación real; no se infiere aprobación
docente ni se construye otro backend ocultamente.

\subsection{Pasos 4--5: contrato y
OpenAPI}\label{informe-pasos-45-contrato-y-openapi}

Se construyen dos archivos OpenAPI 3.1.0 independientes del contrato de
producto:

\begin{enumerate}
\def\labelenumi{\arabic{enumi}.}
\tightlist
\item
  \href{contratos/openapi-observado.yaml}{openapi-observado.yaml}: 14
  operaciones seleccionadas; DTO directos y errores Problem Details,
  parámetros, cookie/CSRF, esquemas y resultados relevantes. Se deriva
  del contrato común y se refina con código y evidencia. No promete
  cobertura de todos los endpoints existentes.
\item
  \href{contratos/openapi-objetivo.yaml}{openapi-objetivo.yaml}:
  \textbf{propuesta no implementada} del sobre
  \path|success/message/data/errors/meta|, conserva códigos y respuestas
  sin cuerpo. Su servidor \path|/api-objetivo-no-implementada| es
  marcador explícito, no una ruta implementada ni una API v2. La
  paginación se mantiene dentro de data como CatalogPage y meta contiene
  correlación; esta variante debe confirmarse, difiere de trasladar
  paginación a meta como alternativa planteada en ACD3.
\end{enumerate}

La selección agrega a las operaciones de catálogo/sesión/proceso la
lectura configurada de modelo. Las rutas siguientes existen en
controladores. Los nombres de parámetros OpenAPI pueden ser
careerId/processId/versionId aunque el controlador use id; representan
la misma posición.

\textbf{Recurso }: Sesión actual

\textbf{Métodos y rutas seleccionadas, bajo /api/v1 }: POST /sessions;
DELETE /sessions

\textbf{Resultado principal }: 201 con SessionView y cookie; 204 sin
cuerpo.

\medskip\hrule\medskip

\textbf{Recurso }: Unidad

\textbf{Métodos y rutas seleccionadas, bajo /api/v1 }: POST
/academic-units

\textbf{Resultado principal }: 201 con unidad, Location y ETag.

\medskip\hrule\medskip

\textbf{Recurso }: Carreras

\textbf{Métodos y rutas seleccionadas, bajo /api/v1 }: GET /careers;
POST /careers

\textbf{Resultado principal }: 200 página visible; 201 alta, Location y
ETag.

\medskip\hrule\medskip

\textbf{Recurso }: Carrera

\textbf{Métodos y rutas seleccionadas, bajo /api/v1 }: GET
/careers/\{careerId\}; PATCH /careers/\{careerId\}

\textbf{Resultado principal }: 200 con recurso y ETag; PATCH exige
CareerInput completo e If-Match.

\medskip\hrule\medskip

\textbf{Recurso }: Historia

\textbf{Métodos y rutas seleccionadas, bajo /api/v1 }: GET
/careers/\{careerId\}/history

\textbf{Resultado principal }: 200 página de revisiones.

\medskip\hrule\medskip

\textbf{Recurso }: Decisión de estado

\textbf{Métodos y rutas seleccionadas, bajo /api/v1 }: POST
/careers/\{careerId\}/state-decisions

\textbf{Resultado principal }: 201 con carrera actualizada; no Location
explícita.

\medskip\hrule\medskip

\textbf{Recurso }: Versión de modelo

\textbf{Métodos y rutas seleccionadas, bajo /api/v1 }: GET
/model-versions/\{versionId\}; GET
/model-versions/\{versionId\}/configured

\textbf{Resultado principal }: DTO legado cuando compatible; 409 si
requiere configuración; ruta configurada devuelve estructura completa.

\medskip\hrule\medskip

\textbf{Recurso }: Procesos

\textbf{Métodos y rutas seleccionadas, bajo /api/v1 }: GET /processes;
POST /processes

\textbf{Resultado principal }: 200 arreglo visible sin paginación; 201
preparación con Location y ETag.

\medskip\hrule\medskip

\textbf{Recurso }: Proceso

\textbf{Métodos y rutas seleccionadas, bajo /api/v1 }: GET
/processes/\{processId\}

\textbf{Resultado principal }: 200 con estado, referencias y ETag.

\medskip\hrule\medskip

Los esquemas precisan campos requeridos, UUID, fechas, límites de
longitud y precondiciones. Reglas no expresables solo por schema:
capacidad/ámbito, unidad activa, código normalizado único, vigencia y
publicación. El contrato académico añade 400 de DTO de catálogo, ausente
en la tabla heredada de respuestas compartidas, y conserva 409 del
modelo legado. No modifica silenciosamente los contratos originales.

Semántica aplicada: GET consulta; POST crea recurso o decisión; PATCH
modifica datos editables del catálogo sin reemplazar ID/estado; DELETE
revoca sesión. No se inventa un PUT de carreras. 409 representa
duplicado/estado incompatible; 412 revisión obsoleta; 428 falta
If-Match; 422 regla de negocio concreta. No devolver 200 con
success=false. El caso 204 carece de schema/cuerpo. El subconjunto
If-Match numérico del proyecto no se enseña como toda la gramática de
HTTP.

\subsection{Paso 6: endpoints base y diferencia respecto de la
guía}\label{informe-paso-6-endpoints-base-y-diferencia-respecto-de-la-guuxeda}

Los endpoints ya estaban implementados en \path|SessionController|,
\path|CareerController|, \path|EvaluationModelController| y
\path|SelfAssessmentController|. Se revisó separación de DTO/HTTP,
aplicación, dominio y repositorio; 13 pruebas ArchUnit comprueban
límites. El servicio de procesos usa directamente
\path|CareerAdministrationService| dentro del mismo módulo; no se
inventa un puerto/microservicio inexistente.

\textbf{No se implementaron nuevos endpoints DRF ni se migró el
producto.} Python/DRF es una exigencia explícita pendiente; el sobre
universal permanece solo especificado. La documentación y los resultados
Spring no bastan para marcar esos pasos como cumplidos. Confirmar
alternativa con docente/equipo y luego implementar lo acordado.

\subsection{Paso 7: verificación
real}\label{informe-paso-7-verificaciuxf3n-real}

Se ejecutaron 23 pruebas existentes seleccionadas con Maven. El intento
inicial dentro del sandbox falló al resolver dependencias; se repitió
con el acceso autorizado y terminó BUILD SUCCESS. Se inició el backend
de clases actuales con \path|sqlite-test|, limitado a 127.0.0.1:18093 y
una base sintética aislada en backend/target. No se usó el JAR antiguo
como si representara los cambios actuales.

Se ejecutó \path|verificacion/verificar_http.py|: primer intento, 15
casos satisfactorios y rechazo inesperado 409 en la representación
legada; se conservó la evidencia. Tras contrastar
\path|LegacyModelVersion.from|/\path|requireLegacy|, se corrigió el
escenario: esperar el rechazo y consultar \path|/configured|. El segundo
recorrido completó \textbf{22 casos}, con estado y aserciones de
cuerpo/cabecera. No se cambió el negocio para forzar el resultado.

Se validaron 22 respuestas capturadas con Ajv contra el contrato
observado; Redocly validó ambos OpenAPI. Swagger UI y
\path|/v3/api-docs| respondieron 200, pero no se ejecutaron operaciones
desde la interfaz Swagger. Los tokens y credenciales de sesión se
redactaron en la evidencia. Los comandos, límites y resultados están en
\url{VERIFICACION.md}.

\subsection{Paso 8 y cierre: trazabilidad, socialización y
versionado}\label{informe-paso-8-y-cierre-trazabilidad-socializaciuxf3n-y-versionado}

La \href{TRAZABILIDAD.md}{matriz APE3} relaciona necesidades,
requisitos, AC, componentes, operaciones y casos reales, diferenciando
cobertura parcial y pendiente. Se integra una
\href{tecnico/APE3.qmd}{página Quarto}, fuentes de diagramas, contratos
y reportes al paquete académico.

Presentación docente, observaciones, aceptación del cambio
tecnológico/sobre y contribuciones individuales siguen pendientes de
registro. La carpeta contiene archivos versionables; \textbf{no
significa que estén comprometidos en Git}. No se creó commit/PR ni se
aplicó GitFlow sobre main. \url{NOTA_INDIVIDUAL.md} distingue trabajo
preparado y aporte verificable que cada persona debe identificar.

\section{Resultados}\label{informe-resultados}

\textbf{Producto esperado }: C4 C1/C2/C3

\textbf{Resultado comprobado }: Structurizr DSL y correspondencia con
ACD3.

\textbf{Límite / pendiente }: Parser/render oficial y validación de
arquitectura con docente.

\medskip\hrule\medskip

\textbf{Producto esperado }: UML/PlantUML

\textbf{Resultado comprobado }: Clases y secuencia vinculadas a código
existente.

\textbf{Límite / pendiente }: Render de diagramas.

\medskip\hrule\medskip

\textbf{Producto esperado }: ADR

\textbf{Resultado comprobado }: Dos propuestas con contexto,
alternativas y consecuencias.

\textbf{Límite / pendiente }: Decisión y aceptación reales.

\medskip\hrule\medskip

\textbf{Producto esperado }: OpenAPI actualizado

\textbf{Resultado comprobado }: Dos documentos académicos de 14
operaciones, lints correctos; 22 muestras conformes al observado.

\textbf{Límite / pendiente }: Propuesta universal no ejecutable; no
sustituye contratos de producto.

\medskip\hrule\medskip

\textbf{Producto esperado }: Endpoints ejecutables

\textbf{Resultado comprobado }: Spring existente arrancado; recorrido
HTTP de catálogo y preparación completado.

\textbf{Límite / pendiente }: No implementados con DRF en esta
actividad.

\medskip\hrule\medskip

\textbf{Producto esperado }: Pruebas

\textbf{Resultado comprobado }: Maven: 23, sin fallos/errores/omitidas;
HTTP: 22; muestras de respuesta: 22.

\textbf{Límite / pendiente }: Sin aceptación PostgreSQL, concurrencia
real multicliente, navegador UI, accesibilidad ni carga.

\medskip\hrule\medskip

\textbf{Producto esperado }: Trazabilidad

\textbf{Resultado comprobado }: Matriz con referencias a AC y casos
exactos.

\textbf{Límite / pendiente }: No certificar requisitos completos por
pruebas parciales.

\medskip\hrule\medskip

\textbf{Producto esperado }: Quarto / informe

\textbf{Resultado comprobado }: Fuente técnica y reporte institucional
preparados.

\textbf{Límite / pendiente }: LaTeX sin compilar; ver evidencia de
render HTML.

\medskip\hrule\medskip

\textbf{Producto esperado }: Versionado / individuos

\textbf{Resultado comprobado }: Archivos, hashes y nota de contribución.

\textbf{Límite / pendiente }: No commit, PR, GitFlow aplicado ni
participación docente acreditada.

\medskip\hrule\medskip

El flujo HTTP cruza cliente, API, servicios y persistencia SQLite;
constituye integración por API, no demostración completa desde Angular
ni validación de migraciones/restricciones PostgreSQL. No se puede
acreditar aún el cierre integral pedido por la guía.

\section{Preguntas de Control}\label{informe-preguntas-de-control}

\subsection{1. ¿Qué diferencia existe entre los niveles C1, C2 y C3 del
modelo
C4?}\label{informe-1-quuxe9-diferencia-existe-entre-los-niveles-c1-c2-y-c3-del-modelo-c4}

C1 ubica personas, sistema y relaciones externas; C2 muestra
aplicaciones/almacenes, responsabilidades y comunicaciones; C3
descompone un contenedor. En SIGEAC: personas↔sistema;
Angular↔Spring↔PostgreSQL; dentro de Spring, API, servicios, dominio,
repositorios, identidad y auditoría. Contenedor C4 no equivale
necesariamente a Docker.

\subsection{2. ¿Por qué un diagrama de clases no debe utilizarse como
sustituto del modelo de
arquitectura?}\label{informe-2-por-quuxe9-un-diagrama-de-clases-no-debe-utilizarse-como-sustituto-del-modelo-de-arquitectura}

Muestra tipos, atributos y dependencias del código, pero no basta para
identificar usuarios, fronteras, despliegues y comunicación. Career y su
repositorio explican el catálogo; no explican por sí solos la relación
del navegador con API/DB. Por eso UML complementa C4.

\subsection{3. ¿Qué relación debe existir entre un requisito, un recurso
y un
endpoint?}\label{informe-3-quuxe9-relaciuxf3n-debe-existir-entre-un-requisito-un-recurso-y-un-endpoint}

El requisito expresa una necesidad comprobable, el recurso representa el
concepto pertinente y el endpoint materializa una operación sobre él.
AA-RF-06 exige registrar carrera; recurso Career; POST /api/v1/careers;
AC-AA-08--12 precisan éxito y rechazos. No derivar necesidades
institucionales solo porque exista una ruta.

\subsection{4. ¿Cómo se selecciona el método HTTP apropiado para una
operación?}\label{informe-4-cuxf3mo-se-selecciona-el-muxe9todo-http-apropiado-para-una-operaciuxf3n}

Por su semántica y efecto deseado, no por el nombre de un método de
programación. El catálogo se consulta con GET y se crea con POST; PATCH
modifica su bloque editable con precondición; DELETE revoca sesión. PUT
implicaría reemplazo y no se inventa aquí. El cliente debe considerar
seguridad semántica e idempotencia al reintentar, sin asumir que POST es
idempotente.

\subsection{5. ¿Qué información mínima debe contener una operación
documentada en
OpenAPI?}\label{informe-5-quuxe9-informaciuxf3n-muxednima-debe-contener-una-operaciuxf3n-documentada-en-openapi}

Ruta y método, propósito/operationId, parámetros con
ubicación/obligatoriedad/schema, requestBody cuando corresponda,
respuestas por estado con media type/schema y cabeceras pertinentes,
seguridad y referencias válidas. En esta práctica también se explicitan
condiciones de negocio y AC; OpenAPI no sustituye toda la
especificación.

\subsection{6. ¿Por qué una decisión arquitectónica relevante debería
documentarse mediante
ADR?}\label{informe-6-por-quuxe9-una-decisiuxf3n-arquitectuxf3nica-relevante-deberuxeda-documentarse-mediante-adr}

Para conservar problema/contexto, restricciones, alternativas, decisión,
estado y consecuencias. Evita presentar una preferencia como acuerdo ya
aprobado. APE3-001 registra la diferencia Spring/DRF; APE3-002, el
impacto de adoptar el sobre universal en clientes v1. Ambos siguen
propuestos.

\subsection{7. ¿Cómo se comprueba que un endpoint conserva trazabilidad
con la
especificación?}\label{informe-7-cuxf3mo-se-comprueba-que-un-endpoint-conserva-trazabilidad-con-la-especificaciuxf3n}

Seguir
necesidad\ensuremath{\rightarrow}RF/HU\ensuremath{\rightarrow}AC\ensuremath{\rightarrow}proceso\ensuremath{\rightarrow}componente\ensuremath{\rightarrow}OpenAPI\ensuremath{\rightarrow}código\ensuremath{\rightarrow}prueba/evidencia
exacta; comprobar resultado y límites. Para AA-RF-08: AC-AA-14/15,
F-ACD3-01, CareerAdministrationService, PATCH career, HTTP-09--12.
Detectar revisión vieja no demuestra por sí solo una carrera real de dos
transacciones ni toda la auditoría: esos límites quedan en la matriz.

\section{Conclusiones}\label{informe-conclusiones}

La arquitectura y el contrato seleccionados tienen correspondencia con
requisitos de Semana 2 y código real. Los 23 tests y 22 intercambios
proporcionan evidencia técnica nueva y limitada del backend Spring; el
rechazo del DTO legado permitió refinar el contrato con una diferencia
observada.

La práctica no está íntegramente acreditada: falta resolver la
tecnología DRF, implementar el sobre aceptado, comprobar PostgreSQL,
renderizar/validar diagramas y presentar el flujo al docente con
evidencia individual/versionada real. Documentar estas diferencias
preserva la honestidad de la entrega.

\section{Recomendaciones}\label{informe-recomendaciones}

\begin{enumerate}
\def\labelenumi{\arabic{enumi}.}
\tightlist
\item
  Resolver APE3-001/002 antes de ampliar código o cambiar DTO públicos.
\item
  Confirmar alcance, actores y retroalimentación; no presentar fixtures
  como normativa validada.
\item
  Repetir campaña y migraciones sobre PostgreSQL, probar permisos de
  lector/ámbito, reglas de vigencia y fallos de auditoría; añadir
  navegador si se exige extremo a extremo web.
\item
  Validar y renderizar DSL/PlantUML; revisar LaTeX antes de entregar.
\item
  Aplicar ramas/revisión/commits acordados y completar aporte por
  integrante con enlaces reales.
\end{enumerate}

\section{Bibliografía /
Referencias}\label{informe-bibliografuxeda--referencias}

Ver \url{FUENTES.md}: requisitos/código/ADR del repositorio y
documentación oficial C4, Structurizr, PlantUML, OpenAPI 3.1.0, RFC
9110/5789/6585/9457, Django REST Framework y Quarto. Las fuentes
técnicas justifican semántica y formatos; no sustituyen decisiones
institucionales o docentes.

\nocite{openapi,structurizr,plantuml,http,drf,quarto}
\bibliographystyle{IEEEtran}
\bibliography{referencias}
\section{Anexos}\label{informe-anexos}

\begin{itemize}
\tightlist
\item
  \path|arquitectura/workspace.dsl|: tres niveles C4.
\item
  \path|arquitectura/catalogo-clases.puml| y
  \path|alta-carrera-secuencia.puml|: UML.
\item
  \path|contratos/openapi-observado.yaml| y
  \path|openapi-objetivo.yaml|.
\item
  \path|evidencias/http-intento-inicial.json|,
  \path|http-resultados.json|, \path|esquemas-resultados.json|,
  \path|pruebas-maven.json| y \path|swagger-disponibilidad.json|.
\item
  \path|VERIFICACION.md|, \path|TRAZABILIDAD.md|,
  \path|NOTA_INDIVIDUAL.md|, \path|OBSERVACIONES.md|,
  \path|tecnico/APE3.qmd| y manifiesto de integridad.
\end{itemize}

\subsection{APE3 --- Trazabilidad de especificación, arquitectura,
contrato e
implementación}\label{traza-ape3--trazabilidad-de-especificaciuxf3n-arquitectura-contrato-e-implementaciuxf3n}

Origen: \href{../AA2-semana-2/DOSSIER_REQUISITOS_AA2.md}{dossier AA2},
\href{../AA2-semana-2/MATRIZ_TRAZABILIDAD.md}{matriz AA2},
\href{../APE2-semana-2/specification.md}{specification APE2} y flujos
F-ACD3-01/02. Los criterios citados mantienen su redacción; una prueba
parcial no los redefine ni los valida por completo. Los RF de esta tabla
son AA-RF y los RNF, AA-RNF.

Evidencia nueva: \path|evidencias/http-resultados.json|,
\path|esquemas-resultados.json|, \path|pruebas-maven.json|. Los HTTP se
refieren a sus campos \path|case|. Estado \textbf{parcial} significa
escenario comprobado en Spring/SQLite, sin certificación de todo el
requisito ni de cumplimiento DRF.

\textbf{Necesidad / actor / historia }: N-AA-01 / A-AA-01 / HU-AA-01

\textbf{Requisito y AC/VC }: AA-RF-01; AC-AA-01/02

\textbf{Proceso y componente }: Acceso; SessionController/SessionService

\textbf{Operación del OpenAPI observado }: POST /sessions

\textbf{Evidencia nueva }: HTTP-02, esquema de SessionView

\textbf{Estado / vacío }: Válido comprobado; credencial errónea/no
enumeración y límite de intentos sin ejecutar.

\medskip\hrule\medskip

\textbf{Necesidad / actor / historia }: N-AA-01 / A-AA-01 / HU-AA-01

\textbf{Requisito y AC/VC }: AA-RF-02; AC-AA-03

\textbf{Proceso y componente }: Acceso; SessionService

\textbf{Operación del OpenAPI observado }: DELETE /sessions y GET
protegido

\textbf{Evidencia nueva }: HTTP-20/21

\textbf{Estado / vacío }: 204 sin cuerpo y acceso revocado comprobados;
no inspección directa del token persistido.

\medskip\hrule\medskip

\textbf{Necesidad / actor / historia }: N-AA-01 / A-AA-01 / HU-AA-01

\textbf{Requisito y AC/VC }: AA-RF-03; AC-AA-04

\textbf{Proceso y componente }: Identidad propia

\textbf{Operación del OpenAPI observado }: GET /me, existente fuera del
subconjunto APE3

\textbf{Evidencia nueva }: Ninguna nueva

\textbf{Estado / vacío }: Pendiente; no ampliar cobertura por login.

\medskip\hrule\medskip

\textbf{Necesidad / actor / historia }: N-AA-01 / A-AA-01--04 /
HU-AA-01--06

\textbf{Requisito y AC/VC }: AA-RF-04; AC-AA-05/06/13/24

\textbf{Proceso y componente }: Autorización por ámbito y CSRF

\textbf{Operación del OpenAPI observado }: Operaciones protegidas

\textbf{Evidencia nueva }: HTTP-01/04/21

\textbf{Estado / vacío }: Falta de sesión y CSRF comprobados; capacidad
funcional, lector ajeno, totales ocultos y proceso ajeno pendientes.

\medskip\hrule\medskip

\textbf{Necesidad / actor / historia }: N-AA-02 / A-AA-02 / HU-AA-02

\textbf{Requisito y AC/VC }: AA-RF-05; AC-AA-07/27

\textbf{Proceso y componente }: F-ACD3-01; CareerAdministrationService

\textbf{Operación del OpenAPI observado }: POST /academic-units

\textbf{Evidencia nueva }: HTTP-03

\textbf{Estado / vacío }: Alta comprobada; no consulta posterior de
unidad ni duplicado de unidad.

\medskip\hrule\medskip

\textbf{Necesidad / actor / historia }: N-AA-02 / A-AA-02 / HU-AA-02

\textbf{Requisito y AC/VC }: AA-RF-06; AC-AA-08--12

\textbf{Proceso y componente }: F-ACD3-01; CareerController/Career

\textbf{Operación del OpenAPI observado }: POST /careers

\textbf{Evidencia nueva }: HTTP-05/06/07

\textbf{Estado / vacío }: Alta, DTO incompleto y duplicado normalizado;
unidad inválida/fechas invertidas pendientes.

\medskip\hrule\medskip

\textbf{Necesidad / actor / historia }: N-AA-04 / A-AA-02/04 / HU-AA-03

\textbf{Requisito y AC/VC }: AA-RF-07; AC-AA-13/14

\textbf{Proceso y componente }: F-ACD3-01; servicio/repositorio

\textbf{Operación del OpenAPI observado }: GET /careers,
/careers/\{careerId\}, /history

\textbf{Evidencia nueva }: HTTP-08/12/13

\textbf{Estado / vacío }: Detalle, lista filtrada e historial; no
aislamiento de lector ni exhaustividad de paginación.

\medskip\hrule\medskip

\textbf{Necesidad / actor / historia }: N-AA-05 / A-AA-02 / HU-AA-04

\textbf{Requisito y AC/VC }: AA-RF-08; AC-AA-14/15

\textbf{Proceso y componente }: F-ACD3-01; servicio y revisiones

\textbf{Operación del OpenAPI observado }: PATCH /careers/\{careerId\}

\textbf{Evidencia nueva }: HTTP-09/10/11/12

\textbf{Estado / vacío }: Revisión nueva, historial y rechazo de
revisión vieja; simulado secuencialmente, no dos transacciones
concurrentes ni ambos actores.

\medskip\hrule\medskip

\textbf{Necesidad / actor / historia }: N-AA-02/03 / A-AA-02/03 /
HU-AA-05/06

\textbf{Requisito y AC/VC }: AA-RF-09; AC-AA-16/17/20

\textbf{Proceso y componente }: F-ACD3-01/02; estado de carrera

\textbf{Operación del OpenAPI observado }: POST
/careers/\{careerId\}/state-decisions

\textbf{Evidencia nueva }: HTTP-15

\textbf{Estado / vacío }: Cambio a DISABLED/revisión 3 comprobado;
bloqueo de proceso posterior y repetición de estado no ejecutados aquí.

\medskip\hrule\medskip

\textbf{Necesidad / actor / historia }: N-AA-03 / A-AA-03 / HU-AA-06

\textbf{Requisito y AC/VC }: AA-RF-10; AC-AA-18/22

\textbf{Proceso y componente }: F-ACD3-02;
EvaluationModelService/ModelCatalogPort

\textbf{Operación del OpenAPI observado }: GET versión legado y
configurada

\textbf{Evidencia nueva }: HTTP-16/16B

\textbf{Estado / vacío }: Identidad y PUBLISHED comprobados en
configurada; 409 legado detectado. Versión no publicada para alta
pendiente.

\medskip\hrule\medskip

\textbf{Necesidad / actor / historia }: N-AA-03 / A-AA-03 / HU-AA-06

\textbf{Requisito y AC/VC }: AA-RF-11; AC-AA-19--22

\textbf{Proceso y componente }: F-ACD3-02; SelfAssessmentService

\textbf{Operación del OpenAPI observado }: POST /processes

\textbf{Evidencia nueva }: HTTP-17

\textbf{Estado / vacío }: PREPARATION/revisión 0 con fixture autorizado;
rechazos de vigencia/modelo/carrera pendientes.

\medskip\hrule\medskip

\textbf{Necesidad / actor / historia }: N-AA-04 / A-AA-03 / HU-AA-06

\textbf{Requisito y AC/VC }: AA-RF-12; AC-AA-23/24

\textbf{Proceso y componente }: F-ACD3-02; servicio de proceso

\textbf{Operación del OpenAPI observado }: GET /processes y
/processes/\{processId\}

\textbf{Evidencia nueva }: HTTP-18/19

\textbf{Estado / vacío }: Consulta propia y presencia en arreglo; actor
ajeno pendiente.

\medskip\hrule\medskip

\textbf{Necesidad / actor / historia }: N-AA-05 / A-AA-02/03 / historias
de cambio

\textbf{Requisito y AC/VC }: AA-RF-13; AC-AA-25

\textbf{Proceso y componente }: Ambos;
AuditPort/SelfAssessmentAuditService

\textbf{Operación del OpenAPI observado }: Mutaciones

\textbf{Evidencia nueva }: Código revisado; sin fallo de auditoría
inducido

\textbf{Estado / vacío }: Pendiente de ensayo de atomicidad/auditoría;
no confundir 201 con cumplimiento del caso de fallo.

\medskip\hrule\medskip

\textbf{Necesidad / actor / historia }: N-AA-05 / A-AA-01--04 /
historias transversales

\textbf{Requisito y AC/VC }: AA-RF-14; AC-AA-26

\textbf{Proceso y componente }: Handlers/CorrelationFilter

\textbf{Operación del OpenAPI observado }: Rechazos de catálogo/modelo

\textbf{Evidencia nueva }: HTTP-01/04/05/07/09/11/14/16; Ajv

\textbf{Estado / vacío }: Cuerpo conforme en muestras; revisión
sistemática de secretos y todas las rutas pendiente.

\medskip\hrule\medskip

\textbf{Necesidad / actor / historia }: N-AA-01 / usuarios

\textbf{Requisito y AC/VC }: AA-RNF-01/02/03; VC-AA-01--03

\textbf{Proceso y componente }: Identidad/acceso

\textbf{Operación del OpenAPI observado }: Sesión y respuestas

\textbf{Evidencia nueva }: Redacción de tokens de evidencia

\textbf{Estado / vacío }: Redactar artefactos no prueba hash,
expiración, privacidad del sistema o ausencia de secretos en logs.

\medskip\hrule\medskip

\textbf{Necesidad / actor / historia }: N-AA-04 / usuarios

\textbf{Requisito y AC/VC }: AA-RNF-04; VC-AA-04

\textbf{Proceso y componente }: Recorridos completos

\textbf{Operación del OpenAPI observado }: Consultas

\textbf{Evidencia nueva }: Duraciones por solicitud en JSON

\textbf{Estado / vacío }: No son p95 visible de navegador, campaña de
carga ni medición del RNF. Pendiente.

\medskip\hrule\medskip

\textbf{Necesidad / actor / historia }: N-AA-05 / A-AA-02/03

\textbf{Requisito y AC/VC }: AA-RNF-05; VC-AA-05

\textbf{Proceso y componente }: Revisiones

\textbf{Operación del OpenAPI observado }: PATCH /careers/\{careerId\}

\textbf{Evidencia nueva }: HTTP-09--12

\textbf{Estado / vacío }: Precondición comprobada; simultaneidad y
actores reales de VC pendientes.

\medskip\hrule\medskip

\textbf{Necesidad / actor / historia }: N-AA-05 / auditor/soporte
autorizado

\textbf{Requisito y AC/VC }: AA-RNF-06/07; VC-AA-06/07

\textbf{Proceso y componente }: Auditoría/correlación

\textbf{Operación del OpenAPI observado }: Transversal

\textbf{Evidencia nueva }: correlationId en errores de muestra

\textbf{Estado / vacío }: No se verificó inmutabilidad ni
correspondencia con logs de éxito/error. Parcial.

\medskip\hrule\medskip

\textbf{Necesidad / actor / historia }: N-AA-01/04 / usuarios

\textbf{Requisito y AC/VC }: AA-RNF-08; VC-AA-08

\textbf{Proceso y componente }: Frontend

\textbf{Operación del OpenAPI observado }: Fuera de prueba HTTP

\textbf{Evidencia nueva }: Ninguna nueva

\textbf{Estado / vacío }: Accesibilidad pendiente, sin inferencia desde
framework.

\medskip\hrule\medskip

\textbf{Necesidad / actor / historia }: N-AA-05 / equipo

\textbf{Requisito y AC/VC }: AA-RNF-09; VC-AA-09

\textbf{Proceso y componente }: Límites C3 y reglas

\textbf{Operación del OpenAPI observado }: Código seleccionado

\textbf{Evidencia nueva }: 13 ModuleArchitectureTest y 10 tests de
catálogo

\textbf{Estado / vacío }: Dependencias y pruebas seleccionadas
correctas; SQL y mantenibilidad completa no certificados.

\medskip\hrule\medskip

\textbf{Necesidad / actor / historia }: N-AA-04/05 / cliente

\textbf{Requisito y AC/VC }: AA-RNF-10; VC-AA-10

\textbf{Proceso y componente }: Contrato/filtros/handlers

\textbf{Operación del OpenAPI observado }: 14 operaciones seleccionadas

\textbf{Evidencia nueva }: Redocly; 22 respuestas Ajv; headers HTTP

\textbf{Estado / vacío }: Muestras comprobadas; sobre objetivo sin
implementación. 401 observado sin WWW-Authenticate: HTTP-AA-01 sigue
abierto.

\medskip\hrule\medskip

La matriz mantiene la cadena
problema/necesidad\ensuremath{\rightarrow}actor\ensuremath{\rightarrow}RF/HU\ensuremath{\rightarrow}AC\ensuremath{\rightarrow}proceso\ensuremath{\rightarrow}C3\ensuremath{\rightarrow}operación\ensuremath{\rightarrow}código\ensuremath{\rightarrow}evidencia.
Los nombres concretos y rutas de implementación se encuentran en
\url{FUENTES.md}, y UML enlaza las clases principales. No se marca
ningún requisito como validado institucionalmente.
\subsection{APE3 --- Verificación reproducible y
resultados}\label{verificacion-ape3--verificaciuxf3n-reproducible-y-resultados}

\subsubsection{Alcance y
condiciones}\label{verificacion-alcance-y-condiciones}

Campaña sobre el backend \textbf{Spring Boot existente}, no DRF. Base
auxiliar sintética y separada:
\path|backend/target/ape3-verification-20261031.sqlite|; puerto local
18093; archivos bajo backend/target. No se usa ni se modifica una base
PostgreSQL institucional. La fecha/hora exactas figuran en los JSON; no
se atribuyen a una sesión de clase.

No se modificó código de aplicación, migraciones ni contratos
originales. La ejecución Maven puede actualizar \path|backend/target|;
la prueba HTTP crea unidades/carreras/procesos sintéticos en la base
aislada. Esa base no se entrega como evidencia normativa ni se versiona.
Las credenciales del perfil demo están ya en fixtures del repositorio;
usar exclusivamente este entorno.

\subsubsection{1. Pruebas existentes
ejecutadas}\label{verificacion-1-pruebas-existentes-ejecutadas}

Desde \path|backend|, se ejecutó:

\begin{Verbatim}[breaklines=true,breakanywhere=true,fontsize=\footnotesize]
.\mvnw.cmd "-Dtest=CareerControllerTest,CareerAdministrationServiceTest,CareerTest,ModuleArchitectureTest" test
\end{Verbatim}

Resultado: \textbf{BUILD SUCCESS; 23 tests, 0 fallos, 0 errores, 0
omitidos}. Desglose: controlador 3; aplicación 5; dominio 2;
arquitectura 13. \href{evidencias/pruebas-maven.json}{Resumen y
nombres}, extraídos de los XML Surefire de esa ejecución. No se copiaron
reportes antiguos de otras suites como si se hubieran ejecutado en APE3.

El primer intento offline falló por parent no disponible para ese
contexto; el intento de red dentro del sandbox falló con
\path|Permission denied: getsockopt|. Tras la autorización de acceso, el
comando terminó correctamente. Estos fallos de entorno no son fallos de
las 23 pruebas ni se ocultan como ejecución exitosa inicial.

MockMvc standalone usa servicio simulado en CareerControllerTest; no
acredita por sí solo filtros, DB o recorrido extremo a extremo. Las
pruebas de aplicación/dominio y ArchUnit tienen límites diferentes.

\subsubsection{2. Arranque temporal
reproducido}\label{verificacion-2-arranque-temporal-reproducido}

Desde \path|backend|:

\begin{Verbatim}[breaklines=true,breakanywhere=true,fontsize=\footnotesize]
.\mvnw.cmd -o "-Dspring-boot.run.useTestClasspath=true" "-Dspring-boot.run.arguments=--spring.profiles.active=sqlite-test --server.address=127.0.0.1 --server.port=18093 --spring.datasource.url=jdbc:sqlite:./target/ape3-verification-20261031.sqlite --sec.files.root=./target/ape3-files" spring-boot:run
\end{Verbatim}

Requiere dependencias previamente disponibles; el arranque final
autorizado usó el contexto que las tenía. Se observó Tomcat en
127.0.0.1:18093 con base \path|/api/v1|, Java 25 y Spring Boot 4.1.1. El
proceso temporal se detuvo al terminar las comprobaciones; no es un
servicio productivo instalado. Para repetir, iniciar de nuevo y detener
con Ctrl+C al finalizar.

Si la ruta de DB ya existe, el cliente genera códigos únicos y no borra
datos. Para ensayos sobre PostgreSQL, preparar un entorno de pruebas
distinto y seguir la configuración principal; no extrapolar garantías de
SQLite.

\subsubsection{3. Recorrido HTTP
ejecutado}\label{verificacion-3-recorrido-http-ejecutado}

Desde la raíz, con el servidor iniciado:

\begin{Verbatim}[breaklines=true,breakanywhere=true,fontsize=\footnotesize]
python documentos-academicos/APE3-semana-3/verificacion/verificar_http.py
\end{Verbatim}

El script exige el host/puerto local indicados y usa urllib, cookie jar,
CSRF y aserciones. No necesita Postman ni Django. Guarda metadatos y
JSON de respuestas redactados; no conserva cookie, contraseña ni
csrfToken auténticos. No muestra tokens en consola.

\textbf{Casos }: HTTP-01/02

\textbf{Comprobación }: Acceso sin sesión 401; login 201

\textbf{Resultado final }: Correcto.

\medskip\hrule\medskip

\textbf{Casos }: HTTP-03

\textbf{Comprobación }: Crear unidad; Location corresponde al ID

\textbf{Resultado final }: 201 correcto.

\medskip\hrule\medskip

\textbf{Casos }: HTTP-04/05

\textbf{Comprobación }: CSRF ausente y DTO incompleto

\textbf{Resultado final }: 403 y 400 correctos.

\medskip\hrule\medskip

\textbf{Casos }: HTTP-06/07

\textbf{Comprobación }: Alta de carrera, ETag 1, Location; duplicado
normalizado

\textbf{Resultado final }: 201 y 409 correctos.

\medskip\hrule\medskip

\textbf{Casos }: HTTP-08--12

\textbf{Comprobación }: Consultar; PATCH sin precondición; cambiar;
revisión obsoleta; historial

\textbf{Resultado final }: 200/428/200/412/200; revisión 2 e historial
de dos elementos.

\medskip\hrule\medskip

\textbf{Casos }: HTTP-13/14

\textbf{Comprobación }: Búsqueda por código e ID inexistente

\textbf{Resultado final }: Página con un resultado; 404.

\medskip\hrule\medskip

\textbf{Casos }: HTTP-15

\textbf{Comprobación }: Registrar desactivación con revisión 2

\textbf{Resultado final }: 201, DISABLED, revisión 3.

\medskip\hrule\medskip

\textbf{Casos }: HTTP-16/16B

\textbf{Comprobación }: Leer modelo sintético configurado mediante DTO
legado / completo

\textbf{Resultado final }: 409 MODEL\_VERSION\_REQUIRES\_V2; 200
PUBLISHED en /configured.

\medskip\hrule\medskip

\textbf{Casos }: HTTP-17--19

\textbf{Comprobación }: Preparar en carrera piloto autorizada; consultar
detalle/lista

\textbf{Resultado final }: 201 PREPARATION/revisión 0/ETag 0; 200 y 200.

\medskip\hrule\medskip

\textbf{Casos }: HTTP-20/21

\textbf{Comprobación }: Logout y acceso posterior

\textbf{Resultado final }: 204, cero bytes; 401 posterior.

\medskip\hrule\medskip

\textbf{Total final: 22 casos completados.}
\href{evidencias/http-resultados.json}{Resultado final}. El
\href{evidencias/http-intento-inicial.json}{intento inicial} conserva la
expectativa errónea de 200 ante el modelo legado y el 409 observado; no
se reescribe para esconder el hallazgo. La corrección fue del escenario,
con fuente \path|LegacyModelVersion.from| /
\path|EvaluationModelService.requireLegacy|, no del código de negocio.

La prueba de concurrencia utiliza revisiones secuenciales dentro de un
cliente. No demuestra dos solicitudes simultáneas, bloqueo de DB ni
resistencia a carga. La prueba no induce fallo de auditoría ni verifica
todos los permisos de lector/ámbito.

\subsubsection{4. Contratos y
muestras}\label{verificacion-4-contratos-y-muestras}

Generación reproducible del subconjunto, desde raíz:

\begin{Verbatim}[breaklines=true,breakanywhere=true,fontsize=\footnotesize]
node documentos-academicos/APE3-semana-3/verificacion/generar_contratos.mjs
node frontend/node_modules/@redocly/cli/bin/cli.js lint documentos-academicos/APE3-semana-3/contratos/openapi-observado.yaml documentos-academicos/APE3-semana-3/contratos/openapi-objetivo.yaml --format=stylish
node documentos-academicos/APE3-semana-3/verificacion/validar_respuestas.mjs
\end{Verbatim}

Se usan dependencias ya instaladas en frontend/node\_modules. Redocly
acepta ambos OpenAPI, cada uno con 14 operaciones. El objetivo declara
explícitamente que no está implementado. Ajv verifica los cuerpos de las
22 muestras contra el observado;
\href{evidencias/esquemas-resultados.json}{resultado}. La comprobación
también exige estado/media type documentados y cuerpo ausente en 204. No
equivale a validar todas las solicitudes, headers, casos negativos,
seguridad ni totalidad de esquemas del producto.

La generación inicial duplicaba CSRF cuando ya estaba por referencia; el
lint lo detectó y se corrigió la detección de parámetros referenciados.
La validación de muestras detectó 400 de DTO ausente en la respuesta de
catálogo heredada; se añadió al subconjunto académico según el handler
real. No se modificaron los contratos originales ni se ocultaron reglas
de lint.

\subsubsection{5. Swagger, diagramas y
documentación}\label{verificacion-5-swagger-diagramas-y-documentaciuxf3n}

GET a \path|/api/v1/swagger-ui/index.html| y \path|/api/v1/v3/api-docs|
devolvió 200 con sus tipos de contenido.
\href{evidencias/swagger-disponibilidad.json}{Evidencia}. Se verificó
disponibilidad HTTP, no interacción visual ni envío de formularios
Swagger. El OpenAPI académico es un artefacto separado del documento
generado por SpringDoc.

Structurizr DSL y PlantUML se contrastaron con código y sintaxis
documentada; sin parser/render oficial local. Quarto renderizó
\path|tecnico/APE3.qmd| a HTML autónomo. Se abrió en Edge mediante
Playwright: tres SVG, sin errores JavaScript ni desbordamiento horizontal
a 1440 px. Se inspeccionó la captura
\path|evidencias/quarto-documentacion.png| y se registró la comprobación
en \path|evidencias/quarto-verificacion.json|. Esta es documentación,
no ejecución de la interfaz Angular. LaTeX se revisa estáticamente,
sin motor ni PDF local.

\subsubsection{6. Límites para la
entrega}\label{verificacion-6-luxedmites-para-la-entrega}

Pendientes: decisión Spring/DRF; implementación del sobre universal;
aceptación PostgreSQL y migraciones; navegador Angular; validaciones
negativas restantes; auditoría/privacidad completas; render de
diagramas; compilación PDF; revisión docente; commits/PR y
contribuciones. El 401 local no incluyó WWW-Authenticate entre los
headers recibidos: aporta evidencia al hallazgo HTTP-AA-01, sin
resolverlo ni introducir Bearer ficticio.
\subsection{APE3 --- Observaciones, decisiones pendientes y
ajustes}\label{obs-ape3--observaciones-decisiones-pendientes-y-ajustes}

\subsubsection{Hallazgos técnicos de esta
preparación}\label{obs-hallazgos-tuxe9cnicos-de-esta-preparaciuxf3n}

\textbf{ID }: APE3-H01

\textbf{Hallazgo / fuente }: Guía DRF frente a Spring/constitución
existentes

\textbf{Tratamiento ya realizado }: APE3-001 propuesto; consulta al
usuario; no migrar aplicación

\textbf{Pendiente }: Elección y acuerdo docente/equipo.

\medskip\hrule\medskip

\textbf{ID }: APE3-H02

\textbf{Hallazgo / fuente }: Sobre universal no implementado

\textbf{Tratamiento ya realizado }: OpenAPI observado y objetivo
diferenciados; APE3-002

\textbf{Pendiente }: Implementación compatible o excepción autorizada.

\medskip\hrule\medskip

\textbf{ID }: APE3-H03

\textbf{Hallazgo / fuente }: Modelo configurado rechaza DTO legado con
409

\textbf{Tratamiento ya realizado }: Conservar intento inicial; verificar
/configured; añadir ruta real al contrato académico

\textbf{Pendiente }: Revisar inventario común de rutas y experiencia
cliente con el equipo.

\medskip\hrule\medskip

\textbf{ID }: APE3-H04

\textbf{Hallazgo / fuente }: Contrato compartido de catálogo no
detallaba 400 DTO

\textbf{Tratamiento ya realizado }: Añadir respuesta real al subconjunto
y validar muestras

\textbf{Pendiente }: Decidir actualización del contrato fuente, fuera de
esta entrega documental.

\medskip\hrule\medskip

\textbf{ID }: APE3-H05

\textbf{Hallazgo / fuente }: Referencia CSRF duplicada al generar
contrato

\textbf{Tratamiento ya realizado }: Corregir generador y repetir lint
sin excepciones

\textbf{Pendiente }: Ninguno respecto de esa duplicación; futuras
regeneraciones requieren lint.

\medskip\hrule\medskip

\textbf{ID }: APE3-H06

\textbf{Hallazgo / fuente }: 401 local sin WWW-Authenticate

\textbf{Tratamiento ya realizado }: Relacionar headers capturados con
HTTP-AA-01

\textbf{Pendiente }: Resolver desafío/protocolo compatible con sesión
actual; no inventar JWT.

\medskip\hrule\medskip

\textbf{ID }: APE3-H07

\textbf{Hallazgo / fuente }: Prueba disponible sin DB PostgreSQL activa

\textbf{Tratamiento ya realizado }: Usar SQLite sintética aislada y
registrar límites

\textbf{Pendiente }: Validar PostgreSQL, migraciones y concurrencia.

\medskip\hrule\medskip

Estas observaciones proceden de inspección/ejecución asistidas;
\textbf{no son retroalimentación docente ni de otros equipos}. Los H de
ACD3 conservan su estado previo; no se reescribe una aceptación
inexistente.

\subsubsection{Registro para la sesión
real}\label{obs-registro-para-la-sesiuxf3n-real}

\textbf{ID }: APE3-OBS-01

\textbf{Fecha / observador real }: Pendiente

\textbf{Observación recibida }: Pendiente de sesión

\textbf{Decisión y motivo }: Pendiente

\textbf{Ajuste o justificación para mantener }: Pendiente

\textbf{Archivo / commit / responsable }: Pendiente de evidencia real

\medskip\hrule\medskip

\subsubsection{Plan para cerrar
pendientes}\label{obs-plan-para-cerrar-pendientes}

\begin{enumerate}
\def\labelenumi{\arabic{enumi}.}
\tightlist
\item
  Presentar alcance y diferencias tecnológicas/contractuales antes de
  afirmar cumplimiento.
\item
  Registrar la alternativa aceptada y actualizar estado de ADR con
  evidencia.
\item
  Si se exige DRF, desarrollar el incremento delimitado y verificarlo;
  no etiquetar la campaña Spring como su sustituto automático.
\item
  Ejecutar el flujo aceptado con base principal y cliente
  correspondiente, registrar errores y pruebas de reglas.
\item
  Renderizar diagramas e informe, incorporar correcciones reales y
  versionar aportes identificables.
\end{enumerate}
\subsection{APE3 --- Contribución individual e historial de
cambios}\label{individual-ape3--contribuciuxf3n-individual-e-historial-de-cambios}

\textbf{Estudiante / equipo / fecha:} pendientes. \textbf{Contribución
personal acreditada:} pendiente de revisión y registro por cada
integrante.

\subsubsection{Cambios respecto de ACD3 y Semana
2}\label{individual-cambios-respecto-de-acd3-y-semana-2}

\begin{itemize}
\tightlist
\item
  Convertir las vistas conceptuales C4 a un workspace Structurizr DSL
  C1/C2/C3 y elaborar UML de clases/secuencia real de catálogo.
\item
  Registrar dos ADR académicos propuestos, sin cambiar decisiones del
  producto.
\item
  Elaborar dos OpenAPI de 14 operaciones: observado y sobre objetivo no
  implementado.
\item
  Refinar respuestas 400 y ruta configurada de modelos por evidencia
  real.
\item
  Ejecutar 23 pruebas existentes, recorrido HTTP final de 22 casos y
  validación de 22 muestras; conservar intento inicial fallido.
\item
  Enlazar RF/AC con componentes, contratos, código y pruebas concretas;
  integrar Quarto y reporte APE.
\end{itemize}

Estos cambios se prepararon con asistencia de IA. Esta lista no
identifica por sí sola a un estudiante como autor humano ni prueba
participación presencial.

\subsubsection{Registro personal para
completar}\label{individual-registro-personal-para-completar}

\textbf{Elemento }: Requisito/criterio revisado

\textbf{Información real del estudiante }: Pendiente.

\medskip\hrule\medskip

\textbf{Elemento }: Fuente y diferencia detectada

\textbf{Información real del estudiante }: Pendiente de ruta/símbolo.

\medskip\hrule\medskip

\textbf{Elemento }: Diagrama/ADR/contrato/código que modificó
personalmente

\textbf{Información real del estudiante }: Pendiente.

\medskip\hrule\medskip

\textbf{Elemento }: Prueba que ejecutó o resultado que verificó
personalmente

\textbf{Información real del estudiante }: Pendiente de caso, fecha y
entorno.

\medskip\hrule\medskip

\textbf{Elemento }: Observación recibida y uso

\textbf{Información real del estudiante }: Pendiente de APE3-OBS real.

\medskip\hrule\medskip

\textbf{Elemento }: Commit / PR / historial verificable

\textbf{Información real del estudiante }: Pendiente; no usar commit
base como aporte APE3.

\medskip\hrule\medskip

\subsubsection{GitFlow: procedimiento propuesto, no
ejecutado}\label{individual-gitflow-procedimiento-propuesto-no-ejecutado}

La rama observada es main y hay trabajo académico no comprometido. Antes
de crear ramas, acordar estrategia con el equipo e identificar archivos
propios/ajenos. Si se usa GitFlow, partir de develop cuando exista y
esté acordada, crear feature/ape3-arquitectura-api, revisar cambios y
preparar commits por contribución real. No afirmar que
develop/feature/PR ya existen ni ejecutar \path|git add .| sobre los
archivos ajenos pendientes.

Seleccionar explícitamente el paquete APE3 y la actualización del índice
académico al preparar el commit; revisar contenido y posibles datos
sensibles. No se hizo commit, push ni PR durante esta preparación. Los
hashes del manifiesto prueban integridad del paquete, no autoría,
GitFlow ni aprobación docente.
\subsection{ADR académico APE3-001 --- Alinear la guía con la base
existente}\label{adr1-adr-acaduxe9mico-ape3-001--alinear-la-guuxeda-con-la-base-existente}

\textbf{Estado: propuesta pendiente de decisión del equipo/docente.} No
sustituye ni cambia ADR-0001/0004/0018 del producto.

\subsubsection{Contexto}\label{adr1-contexto}

La guía requiere Python, Django y Django REST Framework.
\path|backend/pom.xml| configura Java 21/Spring Boot 4.1.1, y
\path|.specify/memory/constitution.md| aprueba
Angular/Java/Spring/PostgreSQL como base del producto.
\path|docs/architecture/research.md| menciona Java 25, mientras la
configuración de compilación es 21; el Java instalado observado es 25.
Son tres datos distintos: planificación, nivel del compilador e
instalación local.

La conversación solicita continuar las entregas académicas del proyecto.
No se ha aportado una aceptación docente para cambiar tecnología ni para
exceptuar el paso Django de APE3. Se consultó al usuario sobre mantener
Spring o añadir un prototipo aislado; no se atribuye una respuesta
inexistente.

\subsubsection{Alternativas}\label{adr1-alternativas}

\textbf{Alternativa }: Documentar y demostrar el incremento Spring
existente

\textbf{Ventaja }: Conserva reglas, seguridad, contratos y cliente ya
presentes

\textbf{Coste / consecuencia }: Requiere autorización docente para
acreditar el paso de Django; el sobre universal tampoco está
implementado.

\medskip\hrule\medskip

\textbf{Alternativa }: Prototipo DRF aislado para el mismo flujo

\textbf{Ventaja }: Permite ejercitar el framework de la guía sin
reemplazar la aplicación

\textbf{Coste / consecuencia }: Hay que delimitar sus requisitos,
autenticación, almacenamiento, tests y relación con el producto; no
atribuirle reglas todavía no portadas.

\medskip\hrule\medskip

\textbf{Alternativa }: Migrar el producto completo a DRF

\textbf{Ventaja }: Unifica tecnología con guía

\textbf{Coste / consecuencia }: Amplía enormemente el alcance, afecta
cliente, datos, seguridad y constitución; no se justifica por una
práctica de tres horas.

\medskip\hrule\medskip

\subsubsection{Tratamiento en este
paquete}\label{adr1-tratamiento-en-este-paquete}

Preparar artefactos verificables del sistema existente y señalar el
cumplimiento Django como \textbf{pendiente}, mientras no exista
elección/validación real. No presentar Java como si fuese DRF ni crear
un backend alternativo sin selección expresa. El contrato actual
conserva DTO directos y Problem Details; el sobre solicitado se define
por separado como objetivo no desplegado.

\subsubsection{Confirmación
necesaria}\label{adr1-confirmaciuxf3n-necesaria}

Registrar alternativa elegida, quién la confirma, fecha, alcance,
consecuencias, evidencia y revisión de requisitos/contrato/pruebas. Si
se elige DRF, deberá implementarse y ejecutarse el incremento acordado;
esta propuesta documental no lo da por completado.
\subsection{ADR académico APE3-002 --- Contrato común y
compatibilidad}\label{adr2-adr-acaduxe9mico-ape3-002--contrato-comuxfan-y-compatibilidad}

\textbf{Estado: propuesta pendiente; sin cambio de API pública.}

\subsubsection{Contexto y
alternativas}\label{adr2-contexto-y-alternativas}

Los controladores actuales retornan Career, AcademicUnit, CatalogPage,
SessionView o listas; los errores son Problem Details. La guía pide
\path|success/message/data/errors/meta|. ADR-0018 preserva DTO v1 y
mantiene separado el contrato objetivo de 005. Cambiar únicamente el
servidor rompería consumidores que leen, por ejemplo, \path|csrfToken|
en raíz o \path|items| de la página.

Alternativas: conservar contrato actual y justificar excepción docente
(compatibilidad inmediata, desviación de guía); migrar de forma
coordinada y versionada al sobre (forma común, coste de
cliente/contratos/tests); prototipo aislado con sobre (evita romper
producto, requiere identificarlo claramente y no suponer cobertura de
negocio completa).

\subsubsection{Propuesta}\label{adr2-propuesta}

Entregar \path|contratos/openapi-observado.yaml| como descripción del
subconjunto existente y \path|contratos/openapi-objetivo.yaml| como
forma académica propuesta, con indicador explícito de no implementación.
No asignar un servidor productivo ficticio ni afirmar disponibilidad de
v2. En el objetivo generado: conservar los códigos HTTP documentados,
colocar representación en
data, errores estructurados en errors y correlación en meta. Las páginas
conservan \path|data.items/page/size/totalElements|; no trasladan ni
duplican paginación en meta. Esta forma coincide con el generador y su
extensión \path|x-pagination-decision|, y sigue pendiente de aceptación
del equipo/docente.
Respuestas 204 no tienen cuerpo; tampoco sobre JSON.

\subsubsection{Consecuencias y decisión
pendiente}\label{adr2-consecuencias-y-decisiuxf3n-pendiente}

Los dos documentos tienen estados distintos y no son intercambiables.
Antes de implementar el objetivo: acordar versión/frontera, actualizar
cliente, filtros, handlers, OpenAPI, pruebas y documentación; comprobar
headers y ausencia de cuerpos. Registrar aprobación real y cambios
versionados. La decisión sigue pendiente y no acredita cumplimiento del
sobre en Spring.
\subsection{APE3 --- Fuentes y localización de
evidencia}\label{fuentes-ape3--fuentes-y-localizaciuxf3n-de-evidencia}

\subsubsection{Base documental y
restricciones}\label{fuentes-base-documental-y-restricciones}

\begin{itemize}
\tightlist
\item
  \path|documentos-academicos/APE2-semana-2/specification.md|: contexto,
  recorte y requisitos locales.
\item
  \path|documentos-academicos/AA2-semana-2/DOSSIER_REQUISITOS_AA2.md|:
  actores, necesidades, RF/RNF, HU, AC y VC.
\item
  \path|documentos-academicos/AA2-semana-2/MATRIZ_TRAZABILIDAD.md|:
  procedencia y límites de validación.
\item
  \path|documentos-academicos/ACD3-semana-3/ACD3_CLINICA_ARQUITECTURA_API.md|:
  F-ACD3-01/02, DA-01--08 y trazabilidad.
\item
  \path|documentos-academicos/ACD3-semana-3/ARQUITECTURA_C4.md|: vistas
  Mermaid de la base existente.
\item
  \path|documentos-academicos/ACD3-semana-3/CONTRATO_API_INICIAL.md|:
  inventario anterior de 18 operaciones y alternativa de sobre.
\item
  \path|documentos-academicos/ACD3-semana-3/REGISTRO_OBSERVACIONES.md|:
  hallazgos técnicos y retroalimentación pendiente, no acta de sesión.
\item
  \path|.specify/memory/constitution.md|: base Spring y límites de
  diseño, pruebas y decisiones abiertas.
\item
  \path|docs/requirements/spec.md| y
  \path|docs/requirements/process-flows.md|: requisitos y procesos de
  producto.
\item
  \path|docs/adr/0001-modular-monolith.md|,
  \path|docs/adr/0004-opaque-sessions-scope-authorization.md|,
  \path|docs/adr/0018-catalog-contract-migration-compatibility.md|:
  decisiones existentes.
\item
  \path|docs/architecture/research.md|: planificación tecnológica; no
  confundir Java 25 allí con nivel 21 configurado en POM.
\end{itemize}

\subsubsection{Implementación y contratos
inspeccionados}\label{fuentes-implementaciuxf3n-y-contratos-inspeccionados}

\textbf{Ruta exacta }: \path|backend/pom.xml|

\textbf{Evidencia utilizada }: Java/Spring, Maven, dependencias y
plugins de prueba/arranque.

\medskip\hrule\medskip

\textbf{Ruta exacta }: \path|backend/src/main/resources/application.yml|

\textbf{Evidencia utilizada }: Base /api/v1, PostgreSQL y configuración
de identidad/archivos.

\medskip\hrule\medskip

\textbf{Ruta exacta
}: \path|backend/src/main/resources/application-sqlite-test.yml|

\textbf{Evidencia utilizada }: Perfil auxiliar, SQL de inicialización y
datos demo.

\medskip\hrule\medskip

\textbf{Ruta exacta
}: \path|backend/src/main/java/ec/edu/unl/sec/identityaccess/api/SessionController.java|

\textbf{Evidencia utilizada }: POST/DELETE sessions; SessionView, cookie
y 204.

\medskip\hrule\medskip

\textbf{Ruta exacta
}: \path|backend/src/main/java/ec/edu/unl/sec/identityaccess/infrastructure/SessionCookie.java|

\textbf{Evidencia utilizada }: SEC\_SESSION, X-CSRF-TOKEN y atributos.

\medskip\hrule\medskip

\textbf{Ruta exacta
}: \path|backend/src/main/java/ec/edu/unl/sec/identityaccess/infrastructure/SessionAuthenticationFilter.java|

\textbf{Evidencia utilizada }: Autenticación y CSRF real antes de
controladores.

\medskip\hrule\medskip

\textbf{Ruta exacta
}: \path|backend/src/main/java/ec/edu/unl/sec/selfassessment/api/CareerController.java|

\textbf{Evidencia utilizada }: Rutas/DTO de catálogo, ETag, Location y
validaciones.

\medskip\hrule\medskip

\textbf{Ruta exacta
}: \path|backend/src/main/java/ec/edu/unl/sec/selfassessment/api/CareerCatalogProblemHandler.java|

\textbf{Evidencia utilizada }: Problem Details de errores de catálogo.

\medskip\hrule\medskip

\textbf{Ruta exacta
}: \path|backend/src/main/java/ec/edu/unl/sec/selfassessment/application/CareerAdministrationService.java|

\textbf{Evidencia utilizada }: Capacidad/ámbito, unicidad, unidad
activa, revisión, auditoría.

\medskip\hrule\medskip

\textbf{Ruta exacta
}: \path|backend/src/main/java/ec/edu/unl/sec/selfassessment/domain/Career.java|

\textbf{Evidencia utilizada }: Data, status, revisión, normalización y
vigencia.

\medskip\hrule\medskip

\textbf{Ruta exacta
}: \path|backend/src/main/java/ec/edu/unl/sec/selfassessment/domain/AcademicUnit.java|

\textbf{Evidencia utilizada }: Record y forma anidada de unidad.

\medskip\hrule\medskip

\textbf{Ruta exacta
}: \path|backend/src/main/java/ec/edu/unl/sec/selfassessment/application/CatalogPage.java|

\textbf{Evidencia utilizada }: items/page/size/totalElements.

\medskip\hrule\medskip

\textbf{Ruta exacta
}: \path|backend/src/main/java/ec/edu/unl/sec/selfassessment/infrastructure/JdbcCareerCatalogRepository.java|

\textbf{Evidencia utilizada }: Adaptador de persistencia/revisiones del
UML.

\medskip\hrule\medskip

\textbf{Ruta exacta
}: \path|backend/src/main/java/ec/edu/unl/sec/selfassessment/api/SelfAssessmentController.java|

\textbf{Evidencia utilizada }: Preparar y consultar proceso, DTO y ETag.

\medskip\hrule\medskip

\textbf{Ruta exacta
}: \path|backend/src/main/java/ec/edu/unl/sec/selfassessment/application/SelfAssessmentService.java|

\textbf{Evidencia utilizada }: Carrera activa, modelo publicado,
PREPARATION/revisión 0.

\medskip\hrule\medskip

\textbf{Ruta exacta
}: \path|backend/src/main/java/ec/edu/unl/sec/selfassessment/domain/SelfAssessmentProcess.java|

\textbf{Evidencia utilizada }: Campos y período de proceso.

\medskip\hrule\medskip

\textbf{Ruta exacta
}: \path|backend/src/main/java/ec/edu/unl/sec/evaluationmodel/api/EvaluationModelController.java|

\textbf{Evidencia utilizada }: Rutas legada/configurada de lectura.

\medskip\hrule\medskip

\textbf{Ruta exacta
}: \path|backend/src/main/java/ec/edu/unl/sec/evaluationmodel/api/LegacyModelVersion.java|

\textbf{Evidencia utilizada }: from llama a requireLegacy.

\medskip\hrule\medskip

\textbf{Ruta exacta
}: \path|backend/src/main/java/ec/edu/unl/sec/evaluationmodel/application/EvaluationModelService.java|

\textbf{Evidencia utilizada }: requireLegacy y rechazo
MODEL\_VERSION\_REQUIRES\_V2.

\medskip\hrule\medskip

\textbf{Ruta exacta
}: \path|backend/src/main/java/ec/edu/unl/sec/selfassessment/infrastructure/SyntheticPilotData.java|

\textbf{Evidencia utilizada }: Identificadores y modelo/carrera piloto
sintéticos.

\medskip\hrule\medskip

\textbf{Ruta exacta
}: \path|backend/src/main/java/ec/edu/unl/sec/selfassessment/infrastructure/SyntheticCareerData.java|

\textbf{Evidencia utilizada }: Carrera/unidad demo y vigencia.

\medskip\hrule\medskip

\textbf{Ruta exacta
}: \path|backend/src/main/java/ec/edu/unl/sec/shared/api/EntityTagSupport.java|

\textbf{Evidencia utilizada }: ETag fuerte numérico y 400/428.

\medskip\hrule\medskip

\textbf{Ruta exacta
}: \path|backend/src/main/java/ec/edu/unl/sec/shared/api/ApiProblemHandler.java|

\textbf{Evidencia utilizada }: 400 de DTO y formato de error.

\medskip\hrule\medskip

\textbf{Ruta exacta }: \path|docs/technical/contracts/openapi.yaml|

\textbf{Evidencia utilizada }: Fuente del subconjunto observado.

\medskip\hrule\medskip

\textbf{Ruta exacta
}: \path|specs/005-complete-evaluation-workflows/contracts/openapi.yaml|

\textbf{Evidencia utilizada }: Referencias de catálogo; no contrato v2
íntegramente implementado.

\medskip\hrule\medskip

\textbf{Ruta exacta }: \path|frontend/package.json|

\textbf{Evidencia utilizada }: Angular y dependencias instaladas de
documentación/validación.

\medskip\hrule\medskip

Los restantes puertos y componentes citados están indexados con ruta
completa en \path|documentos-academicos/ACD3-semana-3/FUENTES.md|. La
lectura no implica haber ejecutado cada componente ni cubierto cada
regla.

\subsubsection{Pruebas ejecutadas y artefactos
nuevos}\label{fuentes-pruebas-ejecutadas-y-artefactos-nuevos}

Se ejecutaron únicamente estas suites de Maven:

\begin{itemize}
\tightlist
\item
  \path|backend/src/test/java/ec/edu/unl/sec/selfassessment/api/CareerControllerTest.java|.
\item
  \path|backend/src/test/java/ec/edu/unl/sec/selfassessment/application/CareerAdministrationServiceTest.java|.
\item
  \path|backend/src/test/java/ec/edu/unl/sec/selfassessment/domain/CareerTest.java|.
\item
  \path|backend/src/test/java/ec/edu/unl/sec/architecture/ModuleArchitectureTest.java|.
\end{itemize}

Sus XML se escribieron en
\path|backend/target/surefire-reports/TEST-ec.edu.unl.sec.…xml|. El
resumen entregado conserva clase, nombres, conteos, tiempo y timestamp.
No se atribuyen a esta actividad las otras suites ya presentes en ese
directorio.

Evidencia HTTP/muestras y herramientas reproducibles:
\path|verificacion/verificar_http.py|, \path|generar_contratos.mjs|,
\path|validar_respuestas.mjs|; resultados bajo \path|evidencias/|.
Plantilla:
\path|.agents/skills/informes-latex-unl/assets/FORMATOAPE.tex|; logo
real: \path|assets/reporte_tecnico_overleaf_logo.png| dentro de la misma
habilidad.

\subsubsection{Referencias técnicas primarias
consultadas}\label{fuentes-referencias-tuxe9cnicas-primarias-consultadas}

\textbf{Fuente }: \href{https://c4model.com/diagrams/system-context}{C4:
contexto}, \href{https://c4model.com/diagrams/container}{contenedores},
\href{https://c4model.com/diagrams/component}{componentes}

\textbf{Relación con el trabajo }: Diferencia entre niveles y frontera
representada; revisadas en ACD3.

\medskip\hrule\medskip

\textbf{Fuente
}: \href{https://docs.structurizr.com/dsl/language}{Structurizr DSL ---
lenguaje} y
\href{https://docs.structurizr.com/dsl/cookbook/component-view/}{vista
de componentes}

\textbf{Relación con el trabajo }: Estructura model/views,
identificadores jerárquicos e inclusiones. No sustituyen un parseo
oficial.

\medskip\hrule\medskip

\textbf{Fuente }: \href{https://plantuml.com/class-diagram}{PlantUML ---
clases}

\textbf{Relación con el trabajo }: Relaciones y tipos reales del
catálogo.

\medskip\hrule\medskip

\textbf{Fuente }: \href{https://spec.openapis.org/oas/v3.1.0}{OpenAPI
3.1.0}

\textbf{Relación con el trabajo }: Paths/operations, schemas, responses
y seguridad; versión elegida compatible con fuente, sin afirmar que es
la más reciente.

\medskip\hrule\medskip

\textbf{Fuente }: \href{https://www.rfc-editor.org/rfc/rfc9110.html}{RFC
9110}

\textbf{Relación con el trabajo }: §9 métodos, §13 precondiciones, §15
estados; ver análisis de ACD3/AA2.

\medskip\hrule\medskip

\textbf{Fuente }: \href{https://www.rfc-editor.org/rfc/rfc5789.html}{RFC
5789}, \href{https://www.rfc-editor.org/rfc/rfc6585.html}{RFC 6585},
\href{https://www.rfc-editor.org/rfc/rfc9457.html}{RFC 9457}

\textbf{Relación con el trabajo }: PATCH, 428/429 y Problem Details;
analizadas en ACD3.

\medskip\hrule\medskip

\textbf{Fuente
}: \href{https://www.django-rest-framework.org/api-guide/views/}{DRF ---
vistas}

\textbf{Relación con el trabajo }: Referencia de tecnología exigida; no
convierte los controladores Java en APIView.

\medskip\hrule\medskip

\textbf{Fuente
}: \href{https://quarto.org/docs/output-formats/html-basics.html}{Quarto
--- HTML}

\textbf{Relación con el trabajo }: Página técnica versionable y salida
HTML.

\medskip\hrule\medskip

Las referencias justifican los formatos/técnicas; las fuentes del
dominio y competencias institucionales siguen en AA2 con sus límites. No
se realizaron entrevistas, validación normativa ni retroalimentación
docente nuevas.

\subsection{Fuente editable: workspace.dsl}
\begin{Verbatim}[breaklines=true,breakanywhere=true,fontsize=\footnotesize]
workspace "SIGEAC APE3" "Catalogo y preparacion de procesos; arquitectura observada, no migracion Django aprobada" {
    !identifiers hierarchical
    model {
        admin = person "Administrador funcional" "Mantener catalogo mediante capacidades autorizadas"
        gestor = person "Gestor autorizado" "Preparar procesos en su ambito"
        lector = person "Lector autorizado" "Consultar informacion visible"
        sigeac = softwareSystem "SIGEAC / SEC-UNL" "Organizar autoevaluacion de carreras con referencias y trazabilidad" {
            web = container "Aplicacion web" "Capturar entradas y presentar resultados" "Angular 22.2.3 / TypeScript"
            api = container "Backend API" "Autorizar y ejecutar casos de uso; monolito modular existente" "Java 21 configurado / Spring Boot 4.1.1" {
                acceso = component "Identidad y acceso" "SessionAuthenticationFilter, SessionService y ScopeAuthorizationService" "Spring Security / sesiones opacas"
                http = component "API de catalogo y procesos" "CareerController y SelfAssessmentController; traducir HTTP" "Spring MVC"
                catalogo = component "Aplicacion de catalogo" "CareerAdministrationService; validar permisos, unidad, unicidad y revision" "Servicio transaccional Java"
                procesos = component "Aplicacion de procesos" "SelfAssessmentService; preparar y consultar sin activar" "Servicio transaccional Java"
                dominio = component "Dominio" "Career, AcademicUnit y SelfAssessmentProcess; invariantes" "Records y reglas Java"
                modelos = component "Catalogo de modelos" "ModelCatalogPort y ModelCatalogAdapter; resolver version publicada" "Puerto y adaptador interno"
                repos = component "Persistencia" "CareerCatalogRepository y SelfAssessmentRepository con adaptadores JDBC" "JDBC / PostgreSQL"
                audit = component "Auditoria" "AuditPort y SelfAssessmentAuditService; registrar cambio critico" "Puerto y servicio Java"
            }
            db = container "Base de datos" "Catalogo, modelos, procesos, identidad y auditoria" "PostgreSQL" {
                tags "Database"
            }
        }
        admin -> sigeac "Mantener catalogo autorizado"
        gestor -> sigeac "Preparar y consultar procesos"
        lector -> sigeac "Consultar datos visibles"
        admin -> sigeac.web "Usar formularios"
        gestor -> sigeac.web "Preparar proceso"
        lector -> sigeac.web "Consultar"
        sigeac.web -> sigeac.api "Invocar API con cookie y CSRF" "HTTP JSON; TLS segun despliegue"
        sigeac.api -> sigeac.db "Persistir y consultar" "SQL / JDBC"
        sigeac.web -> sigeac.api.acceso "Autenticar solicitud" "HTTP"
        sigeac.api.acceso -> sigeac.api.http "Continuar cadena de filtros"
        sigeac.api.http -> sigeac.api.catalogo "Invocar caso de catalogo"
        sigeac.api.http -> sigeac.api.procesos "Invocar caso de proceso"
        sigeac.api.catalogo -> sigeac.api.acceso "Consultar actor y permisos por puertos"
        sigeac.api.procesos -> sigeac.api.acceso "Consultar actor y ambito por puertos"
        sigeac.api.procesos -> sigeac.api.catalogo "Comprobar carrera vigente"
        sigeac.api.procesos -> sigeac.api.modelos "Resolver version publicada"
        sigeac.api.catalogo -> sigeac.api.dominio "Construir y validar"
        sigeac.api.procesos -> sigeac.api.dominio "Construir y validar"
        sigeac.api.catalogo -> sigeac.api.repos "Conservar catalogo e historial"
        sigeac.api.procesos -> sigeac.api.repos "Conservar proceso"
        sigeac.api.catalogo -> sigeac.api.audit "Registrar cambio"
        sigeac.api.procesos -> sigeac.api.audit "Registrar cambio"
        sigeac.api.repos -> sigeac.db "Ejecutar SQL"
        sigeac.api.modelos -> sigeac.db "Consultar mediante infraestructura de modelos"
        sigeac.api.audit -> sigeac.db "Persistir mediante infraestructura de auditoria"
    }
    views {
        systemContext sigeac "C1" {
            include admin gestor lector sigeac
            autolayout lr
        }
        container sigeac "C2" {
            include admin gestor lector sigeac.web sigeac.api sigeac.db
            autolayout lr
        }
        component sigeac.api "C3" {
            include sigeac.web sigeac.db sigeac.api.acceso sigeac.api.http sigeac.api.catalogo sigeac.api.procesos sigeac.api.dominio sigeac.api.modelos sigeac.api.repos sigeac.api.audit
            autolayout lr
        }
        styles {
            element "Person" {
                shape Person
                background #08427b
                color #ffffff
            }
            element "Software System" {
                background #1168bd
                color #ffffff
            }
            element "Container" {
                background #438dd5
                color #ffffff
            }
            element "Database" {
                shape Cylinder
            }
            element "Component" {
                background #85bbf0
                color #000000
            }
        }
    }
}

\end{Verbatim}

\subsection{Fuente editable: catalogo-clases.puml}
\begin{Verbatim}[breaklines=true,breakanywhere=true,fontsize=\footnotesize]
@startuml
title APE3 - Elementos de codigo existentes del catalogo (vista parcial)
skinparam classAttributeIconSize 0
package "api" {
  class CareerController {
    +createCareer(input, request)
    +career(id)
    +updateCareer(id, match, input, request)
    +history(id, page, size)
  }
  class "CareerController.CareerInput" as CareerInput <<record>> {
    code : String
    name : String
    academicUnitId : UUID
    validFrom : LocalDate
    validUntil : LocalDate
  }
}
package "application" {
  class CareerAdministrationService {
    +createCareer(data, correlation) : Career
    +getCareer(id) : Career
    +updateCareer(id, expected, data, correlation) : Career
    +history(id, page, size) : CatalogPage
    +requireActiveCareer(id, start, end)
  }
  interface CareerCatalogRepository
  interface CurrentActorPort
  interface AuthorizationPort
  interface AuditPort
}
package "domain" {
  class Career <<record>> {
    id : UUID
    data : Data
    status : Status
    revision : long
    createdBy : UUID
    createdAt : Instant
    classificationPending : boolean
    +replace(value) : Career
    +state(value) : Career
    +activeOn(date) : boolean
  }
  class "Career.Data" as CareerData <<record>> {
    code : String
    name : String
    academicUnitId : UUID
    validFrom : LocalDate
    validUntil : LocalDate
  }
  enum "Career.Status" as CareerStatus {
    ACTIVE
    DISABLED
  }
}
package "infrastructure" {
  class JdbcCareerCatalogRepository
}
CareerController --> CareerAdministrationService : delegar
CareerController ..> CareerInput : validar DTO
CareerInput ..> CareerData : construir
CareerAdministrationService --> CareerCatalogRepository
CareerAdministrationService --> CurrentActorPort
CareerAdministrationService --> AuthorizationPort
CareerAdministrationService --> AuditPort
CareerAdministrationService ..> Career
Career *-- CareerData
Career --> CareerStatus
JdbcCareerCatalogRepository ..|> CareerCatalogRepository
note bottom of CareerAdministrationService
  Verificar unidad activa, unicidad, permisos y revision.
  API y dominio tienen validaciones adicionales.
  No es una implementacion Django ni un modelo completo.
end note
@enduml

\end{Verbatim}

\subsection{Fuente editable: alta-carrera-secuencia.puml}
\begin{Verbatim}[breaklines=true,breakanywhere=true,fontsize=\footnotesize]
@startuml
title APE3 - Alta de carrera existente; resultado esperado, no captura ejecutada
actor "Administrador funcional" as Admin
participant "CareerController" as API
participant "CareerAdministrationService" as App
participant "CurrentActorPort / AuthorizationPort" as Auth
participant "CareerCatalogRepository" as Repo
participant "AuditPort" as Audit
Admin -> API : POST /api/v1/careers\nCookie + CSRF + JSON
note over API
  Filtros autentican y verifican CSRF antes del controlador.
  DTO invalido puede rechazar antes de llegar al servicio.
end note
API -> App : createCareer(input.data(), correlation)
App -> Auth : requireAdmin()
App -> Repo : findUnit(academicUnitId)
App -> Repo : careerCodeExists(code, null)
alt Unidad activa y codigo disponible
  App -> App : Crear Career ACTIVE, revision 1
  App -> Repo : saveCareer(value, -1, actor, CAREER_CREATED)
  App -> Audit : record(evento)
  App --> API : Career
  API --> Admin : 201 + Location + ETag + Career
else Unidad no disponible o codigo duplicado
  App --> API : CareerCatalogException(422 o 409)
  API --> Admin : Problem Details del handler
end
note over App, Audit
  Servicio transaccional: verificar atomicidad con prueba integrada.
  No inferir commit exitoso solo por este diagrama.
end note
@enduml

\end{Verbatim}
% BEGIN APE3 GUIDE REVIEW
\subsection{APE3 --- Revisión de cumplimiento de la guía
003}\label{revision-guia-ape3--revisiuxf3n-de-cumplimiento-de-la-guuxeda-003}

\subsubsection{Dictamen}\label{revision-guia-dictamen}

\textbf{La entrega tiene desarrollo documental y evidencia técnica
parcial, pero no cumple íntegramente la práctica tal como está
redactada.} Los dos incumplimientos principales son el paso 6 con Django
REST Framework y la aplicación efectiva del sobre universal del paso 4.
El backend demostrado es Spring Boot y devuelve DTO directos/Problem
Details. Una propuesta OpenAPI válida no implementa ese comportamiento
ni concede una excepción docente.

Esta revisión no migra el producto, no acepta ADR, no publica cambios ni
inventa retroalimentación. Se contrastaron los archivos existentes con
la guía aportada, se repitieron validaciones estáticas/de muestras y se
corrigió una inconsistencia documental de paginación.

\textbf{Estados:} Documentado = artefacto disponible; Verificado
parcialmente = comprobación con límites identificados; Pendiente = falta
evidencia o ejecución; No cumplido = hay una diferencia comprobada
respecto de una exigencia expresa. Estos estados no equivalen a
calificación docente.

\subsubsection{1. Datos generales y
objetivos}\label{revision-guia-1-datos-generales-y-objetivos}

\path|APE3_SEMANA_3.md|, tabla inicial, y \path|APE3_SEMANA_3.tex|,
tabla institucional, corresponden a Desarrollo Basado en Plataformas, 5
A, Unidad 1, práctica 003, Edison Leonardo Coronel Romero, Semana 3,
horario por definir, aula/laboratorio y tres horas planificadas. El
resultado de aprendizaje y los cinco objetivos están incorporados.
Integrantes y fecha real siguen pendientes, sin reutilizar nombres de
una plantilla ajena.

El objetivo de arquitectura/trazabilidad tiene desarrollo; el de
contrato tiene propuesta y contraste; el de endpoints aporta ejecución
de implementación previa Spring, \textbf{no nueva implementación DRF}.
No se debe presentar la actividad completa como aprobada.

\subsubsection{2. Procedimiento: evidencia por
actividad}\label{revision-guia-2-procedimiento-evidencia-por-actividad}

\textbf{Punto de la guía }: Inicio 1: revisar Semana 2

\textbf{Evidencia exacta }: Reporte, ``Inicio'';
\path|../APE2-semana-2/specification.md|; dossier/matriz AA2;
\path|TRAZABILIDAD.md|

\textbf{Estado }: Documentado

\textbf{Qué falta / acción concreta }: Confirmar recorte/actores con
equipo; los estados de validación institucional no se infieren.

\medskip\hrule\medskip

\textbf{Punto de la guía }: Inicio 2: seleccionar flujos

\textbf{Evidencia exacta }: Reporte, ``Inicio'' y Paso 1; F-ACD3-01/02

\textbf{Estado }: Documentado

\textbf{Qué falta / acción concreta }: Ratificar que catálogo y
preparación corresponden al incremento de Unidad 1.

\medskip\hrule\medskip

\textbf{Punto de la guía }: Inicio 3: revisar retroalimentación ACD3

\textbf{Evidencia exacta
}: \path|../ACD3-semana-3/REGISTRO_OBSERVACIONES.md|;
\path|OBSERVACIONES.md|

\textbf{Estado }: Pendiente de retroalimentación real

\textbf{Qué falta / acción concreta }: Hay hallazgos técnicos, no
observaciones docentes acreditadas. Incorporar las reales cuando se
reciban.

\medskip\hrule\medskip

\textbf{Punto de la guía }: Paso 1: análisis textual

\textbf{Evidencia exacta }: Reporte, Paso 1; informe ACD3 §2

\textbf{Estado }: Documentado

\textbf{Qué falta / acción concreta }: Incluye participantes,
actividades en infinitivo, decisiones, datos, excepciones y resultados;
validar con actores reales.

\medskip\hrule\medskip

\textbf{Punto de la guía }: Paso 2: C4 y UML

\textbf{Evidencia exacta }: \path|arquitectura/workspace.dsl|, vistas
C1/C2/C3; \path|catalogo-clases.puml|;
\path|alta-carrera-secuencia.puml|

\textbf{Estado }: Documentado; verificación de motores pendiente

\textbf{Qué falta / acción concreta }: Fuentes editables disponibles.
Abrir/validar/renderizar con herramientas correspondientes y corregir si
se detectan errores; los SVG Mermaid del HTML no prueban el parseo de
estos archivos.

\medskip\hrule\medskip

\textbf{Punto de la guía }: Paso 3: ADR

\textbf{Evidencia exacta
}: \path|adr/APE3-001-alineacion-tecnologica.md|;
\path|adr/APE3-002-contrato-universal.md|; ADR-0001/0004/0018 del
producto

\textbf{Estado }: Documentado como propuesta

\textbf{Qué falta / acción concreta }: Elegir y registrar decisiones
reales; no cambiar estado a aceptado sin evidencia.

\medskip\hrule\medskip

\textbf{Punto de la guía }: Paso 4: contrato y sobre universal

\textbf{Evidencia exacta }: \path|contratos/openapi-observado.yaml|;
\path|openapi-objetivo.yaml|; generador

\textbf{Estado }: No cumplido en ejecución del sobre; documentación
disponible

\textbf{Qué falta / acción concreta }: El sobre está solo propuesto.
Implementarlo en la superficie aceptada, incluyendo errores de
filtros/handlers, cliente y exclusiones sin cuerpo; o registrar
excepción docente expresa.

\medskip\hrule\medskip

\textbf{Punto de la guía }: Paso 5: OpenAPI

\textbf{Evidencia exacta }: Dos YAML de 14 operaciones;
\path|evidencias/openapi-lint.json|; revalidación de esta revisión

\textbf{Estado }: Verificado parcialmente

\textbf{Qué falta / acción concreta }: Ambos archivos válidos; observado
y objetivo tienen estados diferentes. Tras implementar el contrato
aceptado, contrastar especificación definitiva y backend.

\medskip\hrule\medskip

\textbf{Punto de la guía }: Paso 6: endpoints DRF

\textbf{Evidencia exacta }: Código actual
\path|backend/src/main/java/ec/edu/unl/sec/|; POM; ADR APE3-001

\textbf{Estado }: No cumplido

\textbf{Qué falta / acción concreta }: Implementación comprobada Spring;
Django/DRF no encontrados en el intérprete Python inspeccionado. Elegir
incremento DRF o justificar excepción aceptada, sin etiquetar Java como
DRF.

\medskip\hrule\medskip

\textbf{Punto de la guía }: Paso 7: solicitudes y validación

\textbf{Evidencia exacta }: \path|verificacion/verificar_http.py|;
\path|evidencias/http-resultados.json|; \path|esquemas-resultados.json|;
\path|pruebas-maven.json|

\textbf{Estado }: Verificado parcialmente

\textbf{Qué falta / acción concreta }: 22 intercambios previos y 23
tests seleccionados sobre Spring/SQLite; se revalidaron solo muestras.
Repetir flujo del backend/contrato aceptados, completar negativos
pertinentes y aceptación PostgreSQL.

\medskip\hrule\medskip

\textbf{Punto de la guía }: Paso 8: trazabilidad/documentación

\textbf{Evidencia exacta }: \path|TRAZABILIDAD.md|;
\path|tecnico/APE3.qmd| y HTML; reporte y fuentes

\textbf{Estado }: Documentado

\textbf{Qué falta / acción concreta }: Conectar futuras
implementaciones/commits con cada AC; no marcar AC completo por
cobertura parcial.

\medskip\hrule\medskip

\textbf{Punto de la guía }: Cierre 1: presentar flujo al docente

\textbf{Evidencia exacta }: Reporte, ``Paso 8 y cierre''; registro vacío
de sesión

\textbf{Estado }: Pendiente

\textbf{Qué falta / acción concreta }: Hay recorrido por cliente
HTTP\ensuremath{\rightarrow}API\ensuremath{\rightarrow}servicios\ensuremath{\rightarrow}SQLite; no hay evidencia de presentación docente. La
guía no exige expresamente Angular: definir el extremo a extremo del
backend seleccionado, sin añadir una obligación de UI no indicada.

\medskip\hrule\medskip

\textbf{Punto de la guía }: Cierre 2: observaciones/ajustes

\textbf{Evidencia exacta }: \path|OBSERVACIONES.md|, APE3-OBS-01

\textbf{Estado }: Pendiente

\textbf{Qué falta / acción concreta }: Registrar persona/fecha,
observación, decisión, ajuste o justificación y evidencia real.

\medskip\hrule\medskip

\textbf{Punto de la guía }: Cierre 3: versionado/aporte

\textbf{Evidencia exacta }: \path|NOTA_INDIVIDUAL.md|; Git observado en
main sin commit de la entrega

\textbf{Estado }: Pendiente

\textbf{Qué falta / acción concreta }: Archivos versionables no son
historial de cambios comprometidos. Aplicar estrategia GitFlow acordada,
revisión y contribuciones identificables; no fabricar commits o autores.

\medskip\hrule\medskip

\subsubsection{3. Herramientas y resultados
esperados}\label{revision-guia-3-herramientas-y-resultados-esperados}

\textbf{Exigencia / resultado }: Git/GitHub y GitFlow

\textbf{Hallazgo verificable }: Existe repositorio; paquete académico no
comprometido; sin evidencia nueva de feature/develop/PR

\textbf{Evaluación }: Falta evidencia del flujo y contribución. No se
publica por el hecho de generar Markdown.

\medskip\hrule\medskip

\textbf{Exigencia / resultado }: Spec Kit

\textbf{Hallazgo verificable }: Constitución, plantillas
\path|.specify/| e insumos SDD presentes

\textbf{Evaluación }: Insumos disponibles; no se acredita nueva
ejecución de comandos Spec Kit durante APE3.

\medskip\hrule\medskip

\textbf{Exigencia / resultado }: Python/Django/DRF

\textbf{Hallazgo verificable }: Python se utilizó como cliente HTTP;
backend es Java

\textbf{Evaluación }: La herramienta cliente no satisface el requisito
de implementar con DRF.

\medskip\hrule\medskip

\textbf{Exigencia / resultado }: PostgreSQL

\textbf{Hallazgo verificable }: Configuración principal existente;
campaña APE3 usó perfil SQLite

\textbf{Evaluación }: No hay aceptación PostgreSQL de esta campaña.

\medskip\hrule\medskip

\textbf{Exigencia / resultado }: C4/UML inicial

\textbf{Hallazgo verificable }: Un DSL con tres vistas y dos fuentes
PlantUML

\textbf{Evaluación }: Artefactos disponibles; render/parseo oficial no
acreditados.

\medskip\hrule\medskip

\textbf{Exigencia / resultado }: Decisiones justificadas

\textbf{Hallazgo verificable }: ADR con contexto, alternativas y
consecuencias

\textbf{Evaluación }: Cumple elaboración documental, sin aceptación
atribuida.

\medskip\hrule\medskip

\textbf{Exigencia / resultado }: OpenAPI/Swagger

\textbf{Hallazgo verificable }: Dos contratos válidos; disponibilidad
HTTP Swagger/api-docs 200 en evidencia previa

\textbf{Evaluación }: No equivale a implementar sobre, ejecutar
operaciones desde Swagger ni validar todos los caminos.

\medskip\hrule\medskip

\textbf{Exigencia / resultado }: Cliente HTTP equivalente

\textbf{Hallazgo verificable }: urllib con cookie jar/CSRF y aserciones
de casos

\textbf{Evaluación }: Es herramienta equivalente admisible; no es
obligatorio repetir con Postman solo por estar instalado.

\medskip\hrule\medskip

\textbf{Exigencia / resultado }: Endpoints implementados/ejecutables

\textbf{Hallazgo verificable }: Spring existente ejecutado, no código
DRF nuevo

\textbf{Evaluación }: Resultado funcional parcial, tecnología exigida
sin cumplir.

\medskip\hrule\medskip

\textbf{Exigencia / resultado }: Respuestas/códigos/validaciones

\textbf{Hallazgo verificable }: Evidencia HTTP y muestras de cuerpo con
Ajv

\textbf{Evaluación }: Cobertura acotada; headers/errores/reglas
restantes y sobre objetivo requieren prueba.

\medskip\hrule\medskip

\textbf{Exigencia / resultado }: Quarto/documentación

\textbf{Hallazgo verificable }: QMD, HTML autónomo y verificación de
documentación con tres SVG

\textbf{Evaluación }: Render y revisión de documentación acreditados; no
confundir con UI del producto.

\medskip\hrule\medskip

\textbf{Exigencia / resultado }: Historial y evidencia individual

\textbf{Hallazgo verificable }: Nota preparada, integrantes/enlaces
pendientes

\textbf{Evaluación }: Falta atribución real y versionado.

\medskip\hrule\medskip

\subsubsection{4. Preguntas de
control}\label{revision-guia-4-preguntas-de-control}

Las siete preguntas están respondidas individualmente en
\path|APE3_SEMANA_3.md|, sección ``Preguntas de Control'', y en el
LaTeX. Se verificó su presencia y correspondencia temática:

\textbf{N.º }: 1

\textbf{Tema }: Diferenciar C1/C2/C3

\textbf{Estado }: Respondida, aplicada a SIGEAC.

\medskip\hrule\medskip

\textbf{N.º }: 2

\textbf{Tema }: Diferenciar clases y arquitectura

\textbf{Estado }: Respondida, UML como complemento.

\medskip\hrule\medskip

\textbf{N.º }: 3

\textbf{Tema }: Relacionar requisito/recurso/endpoint

\textbf{Estado }: Respondida con AA-RF-06 y carrera.

\medskip\hrule\medskip

\textbf{N.º }: 4

\textbf{Tema }: Seleccionar método HTTP

\textbf{Estado }: Respondida sin inventar PUT implementado.

\medskip\hrule\medskip

\textbf{N.º }: 5

\textbf{Tema }: Información mínima OpenAPI

\textbf{Estado }: Respondida con método, ruta, entradas, respuestas,
schemas y seguridad.

\medskip\hrule\medskip

\textbf{N.º }: 6

\textbf{Tema }: Motivo para usar ADR

\textbf{Estado }: Respondida, distingue propuesta y aceptación.

\medskip\hrule\medskip

\textbf{N.º }: 7

\textbf{Tema }: Comprobar trazabilidad

\textbf{Estado }: Respondida con AA-RF-08, AC y casos de revisión;
reconoce límites.

\medskip\hrule\medskip

Tener respuestas correctas documentadas no prueba participación
individual ni compensa los pasos de implementación pendientes.

\subsubsection{5. Correcciones y comprobaciones de esta
revisión}\label{revision-guia-5-correcciones-y-comprobaciones-de-esta-revisiuxf3n}

\begin{enumerate}
\def\labelenumi{\arabic{enumi}.}
\tightlist
\item
  \textbf{Corregida inconsistencia de paginación:} APE3-002 decía
  paginación en meta, pero el objetivo generado conserva CatalogPage
  dentro de data. ADR y anexo LaTeX ahora describen esa forma real de la
  propuesta, mantienen la elección pendiente y no duplican totales. No
  se modificó el contrato ni el negocio.
\item
  \textbf{Diferenciados los niveles de evidencia:} un YAML válido no es
  una API implementada; SVG de documentación no son parseo
  Structurizr/PlantUML; 200 de Swagger no es flujo interactivo; SQLite
  no es aceptación PostgreSQL.
\item
  \textbf{Repetidos lint y validación de muestras:} Redocly acepta ambos
  YAML y Ajv acepta las 22 respuestas guardadas. No se levantó backend
  ni se repitieron Maven o las solicitudes HTTP.
  \path|esquemas-resultados.json| registra la nueva fecha de validación
  de la misma evidencia; los JSON HTTP/Maven originales permanecen
  intactos.
\item
  \textbf{Conservado el manifiesto previo} en
  \path|evidencias/manifest-pre-revision.json|; el nuevo manifiesto
  registra los cambios documentales y límites.
\end{enumerate}

No se resolvieron por documentación las diferencias DRF/sobre, ni se
recibió autorización docente. Las propuestas de AA3 ayudan a discutir
responsabilidad de adaptación, pero no constituyen aprobación ni
implementación en APE3.

\subsubsection{6. Plan concreto para cerrar la
práctica}\label{revision-guia-6-plan-concreto-para-cerrar-la-pruxe1ctica}

\textbf{Prioridad }: Alta

\textbf{Acción }: Resolver diferencia Spring/DRF

\textbf{Evidencia necesaria para cerrarla }: Decisión explícita del
equipo/docente. Si se mantiene la guía, incremento DRF delimitado,
ejecutable y probado; una excepción debe quedar registrada realmente.

\medskip\hrule\medskip

\textbf{Prioridad }: Alta

\textbf{Acción }: Adoptar e implementar el contrato universal acordado

\textbf{Evidencia necesaria para cerrarla }: Backend/cliente
compatibles, política de paginación, manejo de errores y 204, OpenAPI
alineado y pruebas del sobre real.

\medskip\hrule\medskip

\textbf{Prioridad }: Alta

\textbf{Acción }: Demostrar flujo con persistencia principal

\textbf{Evidencia necesaria para cerrarla }: Entorno PostgreSQL de
prueba, migraciones pertinentes y solicitudes de éxito/rechazo con
resultado guardado; sin datos institucionales no autorizados.

\medskip\hrule\medskip

\textbf{Prioridad }: Alta

\textbf{Acción }: Versionar y atribuir aportes

\textbf{Evidencia necesaria para cerrarla }: Ramas/commits/revisión del
flujo acordado, autoría/contribución real y enlaces por integrante.

\medskip\hrule\medskip

\textbf{Prioridad }: Media

\textbf{Acción }: Verificar y visualizar DSL/UML

\textbf{Evidencia necesaria para cerrarla }: Resultado de parser/render
y diagramas revisados, sin presentar Mermaid como captura de otra
herramienta.

\medskip\hrule\medskip

\textbf{Prioridad }: Media

\textbf{Acción }: Completar cobertura de AC del recorte

\textbf{Evidencia necesaria para cerrarla }: Pruebas de ámbitos,
vigencia, modelo no publicado, auditoría y concurrencia cuando
corresponda; actualizar matriz con resultados reales.

\medskip\hrule\medskip

\textbf{Prioridad }: Media

\textbf{Acción }: Realizar cierre docente

\textbf{Evidencia necesaria para cerrarla }: Demostración del flujo,
observaciones recibidas, cambios o justificaciones y responsables.

\medskip\hrule\medskip

\textbf{Prioridad }: Presentación

\textbf{Acción }: Completar nombres/fecha y compilar informe

\textbf{Evidencia necesaria para cerrarla }: Datos reales y revisión de
PDF si se entrega LaTeX; la guía no exige por sí sola PDF, el formato es
apoyo institucional del paquete.

\medskip\hrule\medskip

La falta de un recorrido Angular es un límite adicional de las pruebas
previas, \textbf{no una exigencia independiente de esta guía}. Si el
incremento acordado incluye interfaz, comprobarla; si la demostración
solicitada es de backend, el cliente HTTP puede ser su punto de entrada.

% END APE3 GUIDE REVIEW
\section{Declaración de uso de Inteligencia
Artificial}\label{informe-declaraciuxf3n-de-uso-de-inteligencia-artificial}

Se utilizó ChatGPT/Codex para inspeccionar fuentes, elaborar
arquitectura, contratos, scripts y documentos, y ejecutar
verificaciones. Los resultados declarados provienen de las herramientas;
no se inventaron sesiones de clase, aprobación docente, commits ni
contribuciones humanas. Cada integrante debe revisar el paquete e
identificar su aporte real antes de presentarlo. Esta asistencia no
demuestra participación individual por sí sola.

\end{document}
